\section{Convex norms and Gauduchon degrees}
\label{sec:convex-bundle-norm-obstruction}

For holomorphic line bundles, the trace identity of \cref{prop:gauduchon-trace} becomes a statement about curvature. Let $F$ be a holomorphic line bundle on $X$ and $k$ a smooth Hermitian metric on $F$. In a local holomorphic frame $e$ of $F$ we identify $k$ with the positive function $|e|_k^2$, and the curvature form of $k$ is $R_k=-\ddc\log k$. Let $\Theta$ be a smooth strictly positive real $(n-1,n-1)$-form with $\ddc\Theta=0$; by the root calculation preceding the proof of \cref{thm:lamari}, $\Theta=\gamma^{n-1}$ for a unique Gauduchon metric $\gamma$. The \emph{degree} of $F$ with respect to $\Theta$ is
\[
 \deg_\Theta F=\int_XR_k\wedge\Theta.
\]
Two metrics on $F$ differ by a factor $\mathrm e^{-u}$ with $u$ smooth, which adds $\ddc u$ to the curvature, so the degree is independent of $k$ by \cref{lem:ibp}. It is the Gauduchon degree of the curvature form with respect to $\gamma$; the form $\Theta$ need not be $\mathrm d$-closed.

The squared norm $v\mapsto|v|_k^2$ is plurisubharmonic on the total space of $F$ off the zero section exactly when $\log k$ is plurisubharmonic, that is, when $R_k\leq0$; this follows from the Levi form computation in the proof of \cref{lem:tautological-homogeneous-levi} below. A line bundle of positive degree therefore admits no metric with plurisubharmonic squared norm. This section extends the obstruction to direct sums of line bundles of positive degree, to all holomorphic structures sufficiently close to such a direct sum, and to norms that need not be Hermitian. The mechanism is again integration against $\Theta$: averaging the horizontal complex Hessians of a norm near the line summands produces a negative multiple of the degrees, and this negativity is stable under small changes of the holomorphic structure.

For a smooth complex vector bundle $F$, a \emph{smooth complex norm}\index{Convex norm!complex vector bundle} is a continuous function $p:F\to[0,+\infty)$, positive away from the zero section and smooth there, such that
\[
 p(\lambda v)=|\lambda|p(v),\qquad p(v+w)\leq p(v)+p(w)
\]
for $\lambda\in\mathbb C$ and vectors $v,w$ in the same fiber. Its square $G=p^2$ satisfies $G(\lambda v)=|\lambda|^2G(v)$. Near each line summand of a direct sum, the triangle inequality controls the fiber derivatives of $G$, and these estimates control the averaged horizontal complex Hessians.

\begin{lemma}\label{lem:bundle-degree-normalization}\index{Gauduchon degree!line-bundle normalization}
Fix a smooth positive volume form $\mu$ on $X$. Every holomorphic line bundle $F$ has a Hermitian metric $k$ satisfying
\[
 R_k\wedge\Theta=\lambda\mu,
 \qquad \lambda=\frac{\deg_\Theta F}{\int_X\mu}.
\]
\end{lemma}
\begin{proof}
Start with any metric $k^0$, and put $\mathcal Lu=(\ddc u\wedge\Theta)/\mu$. This scalar operator is uniformly elliptic, has no zeroth-order term, and its kernel consists of the constants by the maximum principle. On H\"older spaces, the homotopy to a Laplace operator used in the proof of \cref{thm:gauduchon} shows that its index is zero, so its cokernel is one-dimensional. Since $\ddc\Theta=0$, integration by parts gives $\int_X\mathcal Lu\,\mu=0$, and therefore its image consists exactly of the smooth functions of zero $\mu$-mean. Hence we can solve
\[
 \ddc u\wedge\Theta=\lambda\mu-R_{k^0}\wedge\Theta,
\]
whose right-hand side has integral zero, and elliptic regularity makes $u$ smooth. Then $k=k^0\mathrm e^{-u}$ has the required curvature.
\end{proof}

Let $F_0=\bigoplus_{s=1}^rF_s$ be a direct sum of holomorphic line bundles, with $r\geq2$. Choose Hermitian metrics $k_s$ on $F_s$ and write $k=\bigoplus_sk_s$. Let $P_s$ be the orthogonal projection onto $F_s$, and let
\[
 \Sigma=\{v\in F_0:|v|_k=1\}
\]
be the unit sphere bundle, with fibers $\Sigma_x$. On each fiber sphere $\Sigma_x$ we use the rotation-invariant probability measure $\mathrm d\sigma_x$. Fix a nonnegative smooth function $\chi$ supported in $(-1,1)$ and positive near zero. For $0<\varepsilon<1/2$, define weights $\rho_s$ on $\Sigma$ by
\[
 \rho_s(v)=c_\varepsilon\chi\left(\frac{|(I-P_s)v|_k^2}{\varepsilon^2}\right),
 \qquad \int_{\Sigma_x}\rho_s\,\mathrm d\sigma_x=1.
\]
In unitary frames respecting the splitting, these weights and the normalizing constant $c_\varepsilon$ are independent of the base point and of the index $s$. Independent phase rotations of the summands preserve them, so they are globally defined. The support of $\rho_s$ lies within distance $\varepsilon$ of the unit circle of $F_s$.

\begin{lemma}\label{lem:convex-norm-axis-estimates}
For a smooth complex norm on $F_0$ with square $G$, put
\[
 a_s(x)=G(e_s),\qquad N(x)=\sum_sa_s(x),\qquad
 D_jG(v)=\mathrm d_vG(P_jv),
\]
where $e_s$ is any $k_s$-unit vector in the fiber of $F_s$ at $x$; complex homogeneity makes $a_s$ well defined. Then
\[
 G(v)\leq N|v|_k^2,\qquad
 |\mathrm d_vG(w)|\leq2N|v|_k|w|_k.
\]
On the support of $\rho_s$,
\[
 |D_jG|\leq2\varepsilon N\quad(j\ne s),\qquad
 |D_sG-2a_s|\leq10\varepsilon N.
\]
\end{lemma}
\begin{proof}
The triangle inequality and the Cauchy--Schwarz inequality give
$p(v)\leq\sum_s|v_s|_k\sqrt{a_s}\leq\sqrt N|v|_k$.
Differentiating $|p(v+tw)-p(v)|\leq|t|p(w)$ in the real variable $t$ gives
$|\mathrm d_vp(w)|\leq p(w)$, and both global bounds follow.

For $v$ in the support of $\rho_s$, set $w=(I-P_s)v$ and $e_s=P_sv/|P_sv|_k$. Then $|w|_k<\varepsilon$ and
$|v-e_s|_k\leq|w|_k+1-|P_sv|_k<2\varepsilon$.
The segment from $e_s$ to $v$ stays in the unit ball and avoids zero, so integrating the derivative bound along it gives $|G(v)-a_s|\leq4\varepsilon N$. For $j\ne s$, we have $|P_jv|_k<\varepsilon$, which gives the first axis bound. Euler's identity $\mathrm d_vG(v)=2G(v)$ yields
\[
 D_sG-2a_s=2(G(v)-a_s)-\sum_{j\ne s}D_jG.
\]
The sum on the right is $\mathrm d_vG(w)$, whose absolute value is at most $2\varepsilon N$. This proves the second axis bound.
\end{proof}

Write $\pi:F_0\to X$ for the bundle projection. For $v\in F_0$ with $\pi(v)=x$, let $H_{0,v}:T_xX\to T_vF_0$ be the horizontal lift for the Chern connection of $k$. It is complex linear, and at unit vectors it is tangent to $\Sigma$, since the Chern connection preserves $k$. For a smooth real function $G$ on $F_0\setminus\{0\}$, define a $(1,1)$-form $\beta_s(G)$ and functions $U_s$, $V_{sj}$ on $X$ by
\[
 \begin{aligned}
 \beta_s(G)&=\int_{\Sigma_x}\rho_s H_{0,v}^*(\ddc G)\,\mathrm d\sigma_x,\\
 U_s&=\int_{\Sigma_x}G\rho_s\,\mathrm d\sigma_x,\\
 V_{sj}&=\int_{\Sigma_x}D_jG\rho_s\,\mathrm d\sigma_x.
 \end{aligned}
\]

\begin{lemma}\label{lem:convex-norm-horizontal-average}\index{Chern connection!horizontal Hessian average}
With $R_j$ the curvature of $k_j$,
\[
 \beta_s(G)=\ddc U_s-\frac12\sum_jR_jV_{sj}.
\]
\end{lemma}
\begin{proof}
At a fixed base point, choose holomorphic frames of the line bundles with
$k_j=\mathrm e^{-\varphi_j}$ and $\mathrm d\varphi_j=0$ at that point. There the horizontal lift has zero vertical component. Write the smooth unitary fiber coordinates as
$\zeta_j=\mathrm e^{-\varphi_j/2}v_j$ and $G(z,v)=g(z,\zeta)$.
At the base point, the first base derivatives of $\zeta_j$ vanish, and its mixed second base derivatives are
$-\frac12(\varphi_j)_{\alpha\bar\beta}\zeta_j$ and their conjugates. The chain rule therefore gives
\[
 G_{\alpha\bar\beta}
 =g_{\alpha\bar\beta}-\frac12\sum_j(\varphi_j)_{\alpha\bar\beta}
       (\zeta_jg_{\zeta_j}+\bar\zeta_jg_{\bar\zeta_j}).
\]
The expression in parentheses is $D_jG$. Both the spherical measure and $\rho_s$ are independent of $z$ in unitary coordinates, so averaging the identity gives the result, since $R_j=\ddc\varphi_j$.
\end{proof}

\begin{theorem}\label{thm:convex-norm-deformation-obstruction}\index{Convex norm!deformation obstruction}\index{Vector bundle!Gauduchon degree obstruction}
Suppose that for each $s$ there is a smooth strictly positive real $(n-1,n-1)$-form $\Theta_s$ satisfying
\[
 \ddc\Theta_s=0,\qquad \deg_{\Theta_s}F_s>0.
\]
Then there is a $C^4$ neighborhood of the split Dolbeault operator on the underlying smooth bundle of $F_0$ such that no integrable Dolbeault operator in this neighborhood admits a smooth complex norm whose square is plurisubharmonic off the zero section.
\end{theorem}
\begin{proof}
\emph{Step 1: the split structure.} Fix $\mu$ and use \cref{lem:bundle-degree-normalization} to choose $k_s$ with
$R_s\wedge\Theta_s=\lambda_s\mu$, where $\lambda_s>0$. Put
\[
 \lambda_* =\min_s\lambda_s,\qquad
 C=5\sum_s\lambda_s+
   \sum_s\sum_{j\ne s}\sup_X\left|\frac{R_j\wedge\Theta_s}{\mu}\right|.
\]
For a squared norm $G$, define $I_0(G)=\sum_s\int_X\beta_s(G)\wedge\Theta_s$. By \cref{lem:convex-norm-horizontal-average}, the term $\ddc U_s$ integrates to zero against $\Theta_s$, and the axis estimates of \cref{lem:convex-norm-axis-estimates} then give
\[
 I_0(G)=-\frac12\sum_{s,j}\int_XV_{sj}R_j\wedge\Theta_s
 \leq-(\lambda_* -\varepsilon C)\int_XN\mu.
\]
Choose $\varepsilon$ so that $\varepsilon C<\lambda_*/2$, and keep all weights fixed.

\emph{Step 2: nearby holomorphic structures.} Write a nearby operator as $\bar\partial_A=\bar\partial_0+A$, with
$\bar\partial_A^2=0$. Its Chern connection for $k$ is $D_A=D_0+A-A^{*k}$, where the adjoint changes the type from $(0,1)$ to $(1,0)$. Let $J_A$ be the corresponding complex structure on the total space and $H_{A,v}$ the horizontal lift of $D_A$. These depend smoothly on $A$, and $H_{A,v}$ is complex linear and tangent to the fixed sphere bundle $\Sigma$.

For $g\in C^\infty(\Sigma,\mathbb R)$, define its homogeneous extension by
\[
 (\mathcal Eg)(v)=|v|_k^2g(v/|v|_k),
\]
and define a real differential operator on $\Sigma$ by
\[
 (\mathcal P_Ag)(v)=\sum_s\rho_s(v)
 \frac{H_{A,v}^*(\ddc_{J_A}\mathcal Eg)_v\wedge\Theta_s(x)}{\mu(x)},
 \qquad x=\pi(v).
\]
It has order at most two. By the identity
$\dc_Jf=-(4\pi)^{-1}\mathrm df\circ J$, its coefficients depend on $A$ and its first derivatives, as well as on the fixed smooth data. Consequently they converge in $C^2$ as $A\to0$ in $C^4$.

\emph{Step 3: the adjoint operator.} Take the formal adjoint $\mathcal P_A^*$ for the fixed density
$\mathrm d\nu=\mathrm d\sigma_x\,\mu(x)$ on the compact sphere bundle, and put $b_A=\mathcal P_A^*1$. Locally, if
$\mathrm d\nu=m(y)\,\mathrm dy$ and
$\mathcal P_A=\sum_{|\alpha|\leq2}c_{\alpha,A}\partial^\alpha$, then
\[
 b_A=m^{-1}\sum_{|\alpha|\leq2}(-1)^{|\alpha|}
                 \partial^\alpha(mc_{\alpha,A}).
\]
Thus $\|b_A-b_0\|_{C^0}\to0$ as $A\to0$ in $C^4$.
For $g=G|_\Sigma$, homogeneity gives $\mathcal Eg=G$ and
\[
 I_A(G):=\int_\Sigma\mathcal P_Ag\,\mathrm d\nu
        =\int_\Sigma gb_A\,\mathrm d\nu.
\]
Since $0<g\leq N\circ\pi$ and the spherical measures have mass one,
\[
 |I_A(G)-I_0(G)|
 \leq\|b_A-b_0\|_{C^0}\int_XN\mu.
\]
Choose the neighborhood so that $\|b_A-b_0\|_{C^0}<\lambda_*/4$ in it. Then
$I_A(G)\leq-\frac14\lambda_*\int_XN\mu<0$ for every smooth squared complex norm $G$. If $G$ were plurisubharmonic for $J_A$, all horizontal pullbacks of $\ddc_{J_A}G$ would be semipositive, and their pairings with $\Theta_s$, weighted by $\rho_s\geq0$, would give $I_A(G)\geq0$, a contradiction.
\end{proof}

The neighborhood depends on the bundle and on the test forms $\Theta_s$, and different line summands may require different tests. Its size is independent of all derivatives of the candidate norms: integration over the fixed sphere bundle transfers those derivatives to the coefficients of the adjoint operator.

The rest of this section develops consequences of \cref{thm:convex-norm-deformation-obstruction}. Norms produced by duality, such as polar norms of Finsler metrics and support functions of neighborhoods of the zero section, are in general only continuous; \cref{subsec:homogeneous-convex-regularization,subsec:tautological-homogeneous-positivity} extend the obstruction to continuous norms with strict convexity margins. Polar duality then shows that, near such a split bundle, the dual of an ample bundle carries no Finsler metric with semipositive Kobayashi curvature (\cref{subsec:polar-finsler-convexity}). Hirzebruch surfaces provide split bundles with positive Gauduchon degrees that are limits of stable ample bundles (\cref{subsec:hirzebruch-ample-comparison,subsec:hirzebruch-stable-extensions,subsec:hirzebruch-stable-ample-deformations}), and a single Stein hypersurface detects the failure of positivity for all Finsler metrics on such a bundle (\cref{subsec:fixed-stein-finsler-hypersurface}). Finally, a sufficiently concave neighborhood of the zero section would give a dual support norm, so the obstruction also bounds the convexity of the complement of the zero section in a projective completion; positive curvature gives bounds in the opposite direction (\cref{subsec:concave-neighbourhood-dual-norms,subsec:projective-complement-curvature-index,subsec:external-tensor-complement-curvature,subsec:flat-restriction-hopf-complements}).

\subsection{Homogeneous convex regularization}
\label{subsec:homogeneous-convex-regularization}
\index{Regularization!homogeneous fibre norms}
\index{Vector bundles!continuous convex norms}
\index{Plurisubharmonic functions!homogeneous regularization}
\index{Greene--Wu gluing}

Let $\pi:F\to X$ be a holomorphic vector bundle over a compact complex manifold. Write $F^\times$ for the complement of its zero section. A continuous real function $Q$ on $F$ is \emph{complex homogeneous of degree two} if
\[
 Q(\lambda v)=|\lambda|^2Q(v).
\]
We call $Q$ \emph{locally strongly convex in the fibres} if, in a local bundle frame and near every point of $F^\times$, there is a constant $c>0$ such that $Q(z,v)-c|v|^2$ is convex on sufficiently small real fibre balls. For continuous functions, strict plurisubharmonicity means that subtracting a sufficiently small positive multiple of the Euclidean squared norm in local total-space coordinates leaves a plurisubharmonic function. Both conditions are local and allow uniform positive margins on compact subsets of $F^\times$.

The smooth-norm obstruction also applies to continuous homogeneous functions with strict positivity margins. The local regularization is Wu's projective-coordinate convolution, written as an average of holomorphic linear shears. Compact gluing follows the argument of Greene and Wu. We give both steps to track strict plurisubharmonicity and strong fibre convexity; the exact attribution is recorded in \cref{sec:history-gauduchon}.

\begin{lemma}\label{lem:compact-sheaf-regularization}
Let $M$ be a compact smooth manifold and let $\mathcal S$ be a sheaf of germs of continuous real functions. Suppose that it is closed under maxima and has the following two properties.

First, if $f$ is a section near a compact set $K$, then $f+u$ is a section near $K$ whenever $u$ is smooth and sufficiently small in the $C^\infty$ topology on a neighbourhood of $K$. Second, every point has a neighbourhood $U$ on which sections admit semilocal smooth approximation: on a compact $K\Subset U$, a section can be approximated uniformly by smooth sections near $K$, with arbitrarily small errors in any prescribed finite collection of derivative norms on compact subsets where the section is already smooth.

Then every global section can be approximated uniformly on $M$ by smooth global sections.
\end{lemma}
\begin{proof}
We give the gluing step. Let $f$ be a global section, smooth near a compact set $L_1$, and let $L_2$ be a compact subset of one of the neighbourhoods $U$ in the statement. Choose relatively compact open sets $W_3,W_4$ with $L_1\subset W_3$, $f$ smooth near $\overline W_3$, and $L_2\subset W_4\Subset U$. Choose $0\leq\rho\leq1$, equal to one near $\overline W_3$, with compact support in the region where $f$ is smooth. Put $V=W_3\cup W_4$. A smooth function $-1\leq\psi\leq1$ can be chosen equal to one near $L_1\cup L_2$ and equal to minus one near $M\setminus V$.

By stability, for sufficiently small $\xi>0$ the function $F=f+\xi\psi$ is still a global section. Approximate $F$ on $\overline W_4$ by a smooth section $g$, with uniform error less than $\delta<\xi/2$. Require also arbitrarily small derivative errors near the compact set $\overline W_4\cap\operatorname{supp}\rho$, where $F$ is smooth. On $V$ define
\[
 h=F\quad\hbox{on }W_3,\qquad
 h=\rho F+(1-\rho)g\quad\hbox{on }W_4.
\]
These definitions agree on the overlap. The function $h$ is smooth. Where $\rho$ vanishes it is the section $g$; where $\rho$ can be nonzero, stability applied to the smooth section $F$ makes $h$ a section if the derivative errors of $g-F$ are sufficiently small. Moreover, $|h-F|<\delta$.

Define $f_1=\max\{f,h\}$ on $V$ and $f_1=f$ outside $V$. Near $L_1\cup L_2$ we have $h>f$, and near the boundary of $V$ we have $h<f$. Thus this formula glues continuously, defines a section by maximum closure, and is smooth near $L_1\cup L_2$. Also $\|f_1-f\|_{C^0}<\xi+\delta$. Near $L_1$, it equals $f+\xi$, so the previous smoothness is preserved.

Take finitely many compact sets $L_2$ of this kind covering $M$, and repeat the step, starting with $L_1$ empty and then taking the union of the previously treated compact sets. At every step choose $\xi$ and $\delta$ small enough that the sum of the finitely many uniform errors is below the prescribed tolerance. The final section is smooth everywhere and has the required uniform error.
\end{proof}

\Needspace{5\baselineskip}
\begin{lemma}\label{lem:homogeneous-shear-regularization}
In a holomorphic bundle chart, use an affine projective coordinate $w=v'/v_r$, where $v=(v',v_r)$ and $v_r\ne0$. Suppose that a positive continuous complex homogeneous function $Q$ is strictly plurisubharmonic and locally strongly convex in the fibres on this cone. Put $q(z,w)=Q(z,(w,1))$. For a nonnegative smooth approximate identity $\eta_\delta(s,t)$ of integral one, define
\[
 q_\delta(z,w)=\int q(z-s,w-t)\eta_\delta(s,t)\,\mathrm ds\,\mathrm dt,
 \qquad Q_\delta(z,v)=|v_r|^2q_\delta(z,v'/v_r).
\]
On every relatively compact projective subchart, sufficiently small $\delta$ gives a smooth positive complex homogeneous function with the same strict plurisubharmonicity and local strong fibre convexity. It converges uniformly to $Q$ on compact normalized cone subsets, and in every finite derivative norm wherever $Q$ is already smooth.
\end{lemma}
\begin{proof}
Complex homogeneity gives the identity
\[
 Q_\delta(z,v)
 =\int Q(z-s,T_tv)\eta_\delta(s,t)\,\mathrm ds\,\mathrm dt,
 \qquad T_t(v',v_r)=(v'-v_rt,v_r).
\]
For fixed $(s,t)$ the map $(z,v)\mapsto(z-s,T_tv)$ is holomorphic and invertible, and its fibre restriction is real linear. On a compact normalized cone subset, all these maps for sufficiently small $(s,t)$ stay in a fixed compact subset of the original cone. Their differentials and inverses are uniformly bounded there. A positive local lower bound for the Levi form of $Q$, interpreted through quadratic barriers, therefore gives a positive lower bound for every pullback and for their average. The same argument with real fibre quadratic barriers preserves local strong convexity. Only fibre segments contained in the local cone neighbourhood are used here.

Ordinary convolution makes $q_\delta$ smooth, and the displayed extension makes $Q_\delta$ smooth on the cone and complex homogeneous. Uniform convergence, derivative convergence on compact smooth regions, and positivity for small $\delta$ follow from the corresponding convolution properties and the positive minimum of $Q$ on compact normalized subsets.
\end{proof}

\begin{theorem}\label{thm:homogeneous-convex-regularization}
Let $Q:F\to[0,+\infty)$ be continuous, positive on $F^\times$, complex homogeneous of degree two, strictly plurisubharmonic on $F^\times$, and locally strongly convex in the fibres. There are continuous complex homogeneous functions $Q_j$, smooth on $F^\times$, strictly plurisubharmonic there and locally strongly convex in the fibres, which converge uniformly to $Q$ on the unit sphere bundle of any fixed smooth Hermitian metric. Each $p_j=\sqrt{Q_j}$ is a smooth complex norm.
\end{theorem}
\begin{proof}
Fix a smooth Hermitian squared norm $G$ and let $M=\mathbb P(F)$ denote the bundle of complex lines. The quotient $m=Q/G$ is a positive continuous function on the compact manifold $M$. Define a sheaf $\mathcal S$ on $M$ by requiring that the corresponding function $Gm$ on the cone be positive, strictly plurisubharmonic, and locally strongly convex in the fibres.

This sheaf is closed under maxima. Indeed, multiplication by $G>0$ commutes with maxima. Near a point of the cone, subtract a common positive quadratic barrier from both functions. Their maximum remains plurisubharmonic; on sufficiently small fibre balls it remains convex. Positivity is also preserved.

For stability on a compact subset $K$ of $M$, restrict to normalized vectors projecting to $K$, and to a fixed compact annulus around them. Finitely many bundle charts provide positive common lower bounds for both barriers. If $u$ is a sufficiently small smooth function on $M$ in $C^2$ near $K$, the Levi form and real fibre Hessian of $Gu$ are smaller than half these bounds on that restricted annulus. The perturbed function is still positive and has both strict margins. Homogeneity transports the local conditions to the rest of the cone over $K$. This proves the stated $C^\infty$ stability.

On a sufficiently small projective coordinate chart, \cref{lem:homogeneous-shear-regularization} gives semilocal approximation, including derivative control on compact subsets where $m$ is smooth. Divide its regularizations by $G$ to obtain local smooth sections of $\mathcal S$. By \cref{lem:compact-sheaf-regularization}, there are global smooth positive sections $m_j\to m$ uniformly. Put $Q_j=Gm_j$ off zero and $Q_j=0$ on the zero section. Homogeneity gives continuity at zero, uniformly over the compact base, and gives the asserted convergence on a fixed unit sphere bundle.

For completeness, local fibre convexity away from zero implies convexity on the whole fibre. The restriction to a line avoiding zero is locally convex and hence convex. On a real line through zero, homogeneity gives $Q_j(tv)=t^2Q_j(v)$ for real $t$, which is convex including at zero. Thus the set $B_x=\{v:Q_j(v)\leq1\}$ is convex, balanced and absorbing. Its gauge is $\sqrt{Q_j}$. If $v,w$ are nonzero, convexity applied to $v/p_j(v)$ and $w/p_j(w)$, with weights proportional to $p_j(v)$ and $p_j(w)$, gives the triangle inequality. The cases with a zero vector follow directly. Positivity, smoothness off zero and complex homogeneity give the other properties of a smooth complex norm.
\end{proof}

\begin{proposition}\label{prop:weak-homogeneous-convex-replacement}
Suppose that $F$ admits a smooth positive complex homogeneous function $H$ of degree two that is strictly plurisubharmonic on $F^\times$. If a positive continuous complex homogeneous function $Q$ of degree two is plurisubharmonic on $F^\times$ and locally strongly convex in the fibres, then $F$ admits a smooth complex norm with strictly plurisubharmonic square on $F^\times$.
\end{proposition}
\begin{proof}
For sufficiently small $\tau>0$, set $Q_\tau=Q+\tau H$. It is positive, complex homogeneous and strictly plurisubharmonic. On a compact normalized annulus, local strong fibre convexity of $Q$ gives a positive uniform margin in finitely many frames. The smooth real fibre Hessian of $H$ is uniformly bounded there. Choose $\tau$ so that its possible negative part consumes less than half the margin. Homogeneity then gives local strong fibre convexity everywhere off zero. Apply \cref{thm:homogeneous-convex-regularization} to $Q_\tau$ and take one of its smooth regularizations. Its square root is the required smooth complex norm.
\end{proof}

\begin{corollary}\label{cor:continuous-convex-norm-obstruction}
For every holomorphic structure in the neighbourhood of \cref{thm:convex-norm-deformation-obstruction}, there is no positive continuous complex homogeneous function of degree two that is both strictly plurisubharmonic off zero and locally strongly convex in the fibres. If that structure admits a function $H$ as in \cref{prop:weak-homogeneous-convex-replacement}, there is also no positive continuous complex homogeneous function of degree two that is plurisubharmonic off zero and locally strongly convex in the fibres.
\end{corollary}
\begin{proof}
The first assertion follows from \cref{thm:homogeneous-convex-regularization}, which would produce a smooth complex norm with plurisubharmonic square, contrary to \cref{thm:convex-norm-deformation-obstruction}. For the second assertion use \cref{prop:weak-homogeneous-convex-replacement} and the same obstruction.
\end{proof}

The second assertion requires the specified strictly plurisubharmonic homogeneous function $H$. The argument does not assert that this function exists for every holomorphic bundle or follows from any positivity hypothesis on the base. The regularization does not give a deformation obstruction for continuous fibre norms without a strong convexity margin.

\subsection{Strict positivity from the tautological metric}
\label{subsec:tautological-homogeneous-positivity}
\index{Vector bundle!tautological positivity}
\index{Tautological line bundle!homogeneous Levi form}

The second part of \cref{cor:continuous-convex-norm-obstruction} requires a smooth strictly plurisubharmonic function on $F^\times$ that is complex homogeneous of degree two. Such a function is the squared norm of a vector, viewed as a point of the tautological line it spans, for a suitable metric on the tautological line bundle. Its Levi form is strictly positive exactly when the dual tautological metric is positively curved, so the function exists when $F^*$ is ample.

Let $F\to X$ be a holomorphic vector bundle of positive rank. In this subsection, $\mathbb P(F)$ parametrizes complex lines in $F$, $\mathcal O_{\mathbb P(F)}(-1)$ is the tautological line bundle whose fiber over a line is that line, and
\[
 L=\mathcal O_{\mathbb P(F)}(1)
   =\mathcal O_{\mathbb P(F)}(-1)^*.
\]
The map $v\mapsto([v],v)$ identifies $F^\times$ holomorphically with the complement of the zero section in $\mathcal O_{\mathbb P(F)}(-1)$. We use the curvature convention $R_h=-\ddc\log h$ of \cref{sec:convex-bundle-norm-obstruction}.

\Needspace{5\baselineskip}
\begin{lemma}\label{lem:tautological-homogeneous-levi}
A smooth Hermitian metric $h$ on $L$ determines a positive smooth complex homogeneous function $H$ of degree two on $F^\times$: the squared norm of $v$ in the dual metric on the tautological line $\mathcal O_{\mathbb P(F)}(-1)_{[v]}$. The function $H$ is strictly plurisubharmonic if and only if $R_h$ is strictly positive on $\mathbb P(F)$.
\end{lemma}
\begin{proof}
Take local holomorphic coordinates $y$ on $\mathbb P(F)$ and a holomorphic frame $e$ of the tautological line bundle, and write the metric of the dual frame as $h=\mathrm e^{-\varphi}$. The dual metric gives
\[
 H(y,\lambda e)=|\lambda|^2\mathrm e^{\varphi(y)},\qquad
 R_h=\ddc\varphi.
\]
Let $\mathcal L$ denote the complex Hessian, without the factor $\mathrm i/(2\pi)$ of the curvature normalization. Differentiating at a point with $\lambda\ne0$ gives
\[
 \mathcal L H(u,t)
 =H\left(\mathcal L\varphi(u)+
          \left|\frac{t}{\lambda}+\partial\varphi(u)\right|^2\right),
 \qquad (u,t)\in T_y^{1,0}\mathbb P(F)\oplus\mathbb C.
\]
If $\mathcal L\varphi$ is positive definite, the right-hand side is positive for every nonzero $(u,t)$: the first term detects $u\ne0$, and the second detects $t\ne0$ when $u=0$. Conversely, for $u\ne0$, set $t=-\lambda\partial\varphi(u)$; strict positivity of $\mathcal L H$ then forces $\mathcal L\varphi(u)>0$. This is exactly positivity of $R_h$. The squared tautological norm is independent of the chosen frame and is complex homogeneous, so the local functions define the asserted global function $H$.
\end{proof}

A line bundle is very ample when its complete linear system defines a holomorphic embedding, and ample when some positive tensor power is very ample. With the convention above, a vector bundle $F^*$ is called ample when $L$ is ample.

\Needspace{5\baselineskip}
\begin{proposition}\label{prop:ample-dual-strict-homogeneous-function}
\index{Ampleness!strict homogeneous functions on the dual}
If $F^*$ is ample, there is a positive complex homogeneous function $H$ of degree two on $F^\times$ that is smooth and strictly plurisubharmonic there and, after extension by zero, continuous on $F$.
\end{proposition}
\begin{proof}
Choose $m\geq1$ such that $L^m$ is very ample. Its complete linear system gives an embedding $\iota\colon\mathbb P(F)\to\mathbb P^N$ with $L^m=\iota^*\mathcal O_{\mathbb P^N}(1)$. Pull back the Fubini--Study metric to a metric $k$ on $L^m$, and take its positive $m$th root $h$ on $L$: in a local frame $h=k^{1/m}$, and the transformation law for $L^m$ makes these local roots a metric on $L$. Its curvature is
\[
 R_h=\frac1m\iota^*\omega_{\mathrm{FS}}>0,
\]
where $\omega_{\mathrm{FS}}$ is the curvature form of the Fubini--Study metric on $\mathcal O_{\mathbb P^N}(1)$; the pullback is strictly positive because the differential of the embedding is injective. Apply \cref{lem:tautological-homogeneous-levi}. Since the base is compact, the resulting $H$ is bounded on the unit sphere bundle of a fixed Hermitian metric. Homogeneity then gives $0<H(v)\leq C|v|^2$ off zero, which proves continuity of the extension by zero.
\end{proof}

\Needspace{5\baselineskip}
\begin{corollary}\label{cor:ample-dual-continuous-convex-obstruction}
Let $F$ carry a holomorphic structure in the neighborhood of \cref{thm:convex-norm-deformation-obstruction}. If $F^*$ is ample, there is no positive continuous complex homogeneous function of degree two on $F$ that is plurisubharmonic off zero and locally strongly convex in the fibers.
\end{corollary}
\begin{proof}
By \cref{prop:ample-dual-strict-homogeneous-function}, ampleness of $F^*$ gives the function $H$ required in \cref{prop:weak-homogeneous-convex-replacement}, so the second assertion of \cref{cor:continuous-convex-norm-obstruction} applies.
\end{proof}

The function $H$ is not the square of a Hermitian norm in general, since the metric on the tautological line may vary with the line. It provides the strict plurisubharmonicity on the total space needed for regularization, whereas local strong convexity in the fibers remains a separate hypothesis in the continuous obstruction.

\subsection{Polar norms of semipositive Finsler metrics}
\label{subsec:polar-finsler-convexity}
\index{Kobayashi curvature!polar norm}
\index{Finsler metric!semipositive curvature}

Let $E\to X$ be a holomorphic vector bundle of positive rank over a compact complex manifold of dimension $n\geq1$. Write $E^\times$ for the complement of its zero section. A squared Finsler metric is a positive smooth function $G$ on $E^\times$ with
\[
 G(x,tv)=|t|^2G(x,v),\qquad t\in\mathbb C\setminus\{0\},
\]
extended continuously by zero. Strong pseudoconvexity means that its vertical complex Hessian $(G_{i\bar j})$ is positive definite. It does not require convexity in the real fibre variables.

In local holomorphic base and fibre coordinates $(z,v)$, put
\[
 \begin{gathered}
 S_{\alpha\bar\beta}
 =G_{\alpha\bar\beta}
  -G_{\alpha\bar j}G^{\bar j i}G_{i\bar\beta},\\
 \mathcal K_G=-\frac{\mathrm i}{2\pi G}
 S_{\alpha\bar\beta}\,\mathrm dz^\alpha\wedge\mathrm d\bar z^\beta,
 \end{gathered}
\]
where $(G^{\bar j i})$ is the inverse vertical Hessian. This is the Kobayashi curvature convention used here. Semipositivity of $\mathcal K_G$ means $S\leq0$.

\Needspace{6\baselineskip}
\begin{lemma}\label{lem:polar-finsler-osculating-section}
\index{Holomorphic section!osculating quadratic jet}
If $G$ is strongly pseudoconvex and $\mathcal K_G>0$, then for every $v_0\in E_{x_0}^\times$ there is a local holomorphic section $\sigma$ with $\sigma(x_0)=v_0$ and
\[
 G(z,\sigma(z))\leq G(x_0,v_0).
\]
For $\varphi=|z|^2$ and $\varepsilon>0$, the metric $G_\varepsilon=\mathrm e^{-\varepsilon\varphi}G$ satisfies
\[
 \mathcal K_{G_\varepsilon}=\mathcal K_G+\varepsilon\ddc\varphi.
\]
\end{lemma}
\begin{proof}
The vertical Hessian is positive, so its Levi-orthogonal complement $H_{G,v}$ projects isomorphically onto the base tangent space. Its Levi form is the Schur complement $S$, and hence is negative definite when $\mathcal K_G>0$. For the vertical radial vector $R=v_i\partial_{v_i}$, complex homogeneity and differentiation give
\[
 \partial G(R)=G,\qquad
 \mathcal L_G(a,\bar R)=\partial G(a).
\]
Since $R$ is vertical, the second identity implies $\partial G|_{H_{G,v}}=0$.

Choose a holomorphic section $\sigma_0$ through $v_0$ with first jet equal to $H_{G,v_0}$. In coordinates centred at $x_0$, the function $G\circ\sigma_0$ has zero first derivative and Taylor expansion
\[
 G(x_0,v_0)+P(z)+\overline{P(z)}
 +S_{\alpha\bar\beta}z^\alpha\bar z^\beta+O(|z|^3),
\]
where $P$ is a homogeneous holomorphic quadratic polynomial. The identity $\partial_vG(v_0)(v_0)=G(v_0)>0$ gives a vertical vector $b$ with $\partial_vG(v_0)(b)=1$. Set $\sigma=\sigma_0-bP$ in the local frame. This changes only the pure quadratic terms, cancelling $P+\bar P$. The remaining quadratic form is negative definite, and therefore dominates the cubic remainder on a sufficiently small ball. The section stays nonzero there and has the asserted inequality.

For the last assertion, differentiate $\mathrm e^{-\varepsilon\varphi}G$ in the displayed Schur complement. Euler's identity gives $G_{\bar j}=v_iG_{i\bar j}$ and $G_\alpha=G_{\alpha\bar j}\bar v_j$. Thus the terms containing $\partial\varphi$, including its square, cancel between the base Hessian and the mixed block product, leaving
\[
 S(G_\varepsilon)=\mathrm e^{-\varepsilon\varphi}
       \bigl(S(G)-\varepsilon G(\varphi_{\alpha\bar\beta})\bigr).
\]
Division by $G_\varepsilon$ proves the curvature formula. Its vertical Hessian is the positive multiple $\mathrm e^{-\varepsilon\varphi}(G_{i\bar j})$.
\end{proof}

\Needspace{6\baselineskip}
\begin{proposition}\label{prop:semipositive-finsler-polar-psh}
\index{Polar norm!plurisubharmonic square}
Suppose $G$ is strongly pseudoconvex and $\mathcal K_G\geq0$. Its polar square
\[
 Q(x,\xi)=\max_{[v]\in\mathbb P(E_x)}
                \frac{|\xi(v)|^2}{G(x,v)},\qquad \xi\in E_x^*,
\]
is continuous and plurisubharmonic on $E^*$, positive off zero and complex homogeneous of degree two. The function $\log Q$ is plurisubharmonic off zero, and $\sqrt Q$ is a complex norm in each fibre.
\end{proposition}
\begin{proof}
The ratio is invariant under complex rescaling of $v$. Maximization on the compact unit sphere in a local frame proves continuity. Positive upper and lower bounds for $G$ on that sphere compare $Q$ with a Hermitian squared dual norm. They prove positivity off zero, continuity at zero and complex homogeneity. Moreover,
\[
 \sqrt{Q(x,\xi)}=\sup_{v\ne0}\frac{|\xi(v)|}{\sqrt{G(x,v)}}.
\]
Each term satisfies the triangle inequality in $\xi$, so their supremum does too.

On a coordinate ball set $G_\varepsilon=\mathrm e^{-\varepsilon|z|^2}G$. By \cref{lem:polar-finsler-osculating-section}, its curvature is strictly positive and its polar square is $Q_\varepsilon=\mathrm e^{\varepsilon|z|^2}Q$. At $(x_0,\xi_0)$ with $\xi_0\ne0$, choose a maximizing line and $v_0$ on it. The lemma supplies a holomorphic section $\sigma$ such that
\[
 \log Q_\varepsilon(z,\xi)
 \geq\log|\xi(\sigma(z))|^2-\log G_\varepsilon(x_0,v_0),
\]
with equality at $(x_0,\xi_0)$. The pairing is nonzero on a neighbourhood of that point, so the right side is pluriharmonic there. Restriction to any sufficiently small holomorphic disc centred at the point therefore gives the submean inequality for $\log Q_\varepsilon$. This argument applies at every point off zero; continuity gives plurisubharmonicity there.

As $\varepsilon\downarrow0$, $\log Q_\varepsilon=\log Q+\varepsilon|z|^2$ converges locally uniformly off zero. Passing to the disc submean inequalities proves the assertion for $\log Q$. The exponential is increasing and convex, so Jensen's inequality makes $Q$ plurisubharmonic off zero. At a point of the zero section, $Q=0$ and $Q\geq0$; hence every centred disc already satisfies its submean inequality there. Together with continuity and the local disc inequalities at nonzero points, this proves plurisubharmonicity on all of $E^*$.
\end{proof}

\Needspace{6\baselineskip}
\begin{lemma}\label{lem:polar-finsler-uniform-strong-convexity}
\index{Legendre duality!strong fibre convexity}
Let $k$ be a smooth Hermitian metric on $E$. For every positive smooth complex homogeneous $G$ of degree two on $E^\times$, its polar square $Q$ has a constant $c>0$ such that
\[
 Q(x,\xi)-c|\xi|_{k^*}^2
\]
is real convex on every dual fibre. Strong pseudoconvexity and real convexity of $G$ are not required for this assertion.
\end{lemma}
\begin{proof}
The real fibre Hessian of $G$ is bounded above by $2L$ times $k$ on its compact unit sphere bundle, for some $L>0$. Homogeneity extends this bound to every nonzero vector. Also $G$ is $C^1$ at zero with derivative zero: its value is $O(|v|^2)$ and its first fibre derivative is $O(|v|)$, uniformly over the base. Integrating the Hessian estimate along a segment gives
\begin{equation}\label{eq:polar-finsler-upper-quadratic}
 G(v+w)\leq G(v)+\mathrm d_vG(w)+L|w|_k^2.
\end{equation}
A segment passing through zero is handled by splitting there and using continuity of the first derivative.

Complex homogeneity and completing the square on each complex line give the exact identity
\begin{equation}\label{eq:polar-finsler-legendre}
 Q(x,\xi)=\max_{v\in E_x}
                  \{2\operatorname{Re}\xi(v)-G(x,v)\}.
\end{equation}
The maximum exists because $G(v)\geq a|v|_k^2$ for a uniform $a>0$. Fix a maximizer $v_0$ for $\xi_0$. Differentiation, including when $\xi_0=v_0=0$, gives $\mathrm dG(v_0)=2\operatorname{Re}\xi_0$. Evaluate \eqref{eq:polar-finsler-legendre} for $\xi_0+\zeta$ at $v_0+w$ and apply \eqref{eq:polar-finsler-upper-quadratic}. The terms involving $\xi_0(w)$ cancel, giving
\[
 Q(\xi_0+\zeta)\geq Q(\xi_0)+2\operatorname{Re}\zeta(v_0)
                 +2\operatorname{Re}\zeta(w)-L|w|_k^2.
\]
Taking the supremum in $w$ yields
\[
 Q(\xi_0+\zeta)\geq Q(\xi_0)+2\operatorname{Re}\zeta(v_0)
                         +L^{-1}|\zeta|_{k^*}^2.
\]
Set $c=L^{-1}$ and $\Phi=Q-c|\cdot|_{k^*}^2$. After subtracting the square, this inequality gives a global real affine support below $\Phi$, equal to $\Phi$ at $\xi_0$. The supremum of these supports is exactly $\Phi$, so $\Phi$ is convex. Neither uniqueness of the maximizer nor real convexity of $G$ is used.
\end{proof}

\Needspace{6\baselineskip}
\begin{corollary}\label{cor:ample-semipositive-finsler-smooth-dual}
If $E$ is ample and has a strongly pseudoconvex squared Finsler metric with semipositive Kobayashi curvature, then $E^*$ admits a smooth complex norm whose square is strictly plurisubharmonic off zero and locally strongly convex in the real fibres.
\end{corollary}
\begin{proof}
The polar square is a positive continuous plurisubharmonic complex homogeneous function by \cref{prop:semipositive-finsler-polar-psh}. It has the strong fibre convexity required in \cref{prop:weak-homogeneous-convex-replacement}, by \cref{lem:polar-finsler-uniform-strong-convexity}. Apply \cref{prop:ample-dual-strict-homogeneous-function} to $F=E^*$, since $F^*=E$ is ample. It supplies the positive smooth strictly plurisubharmonic homogeneous function $H$ on $E^*$. The construction in \cref{prop:weak-homogeneous-convex-replacement} uses \cref{thm:homogeneous-convex-regularization}, which retains strong fibre convexity. Its square root is the asserted smooth norm.
\end{proof}

\Needspace{6\baselineskip}
\begin{corollary}\label{cor:gauduchon-semipositive-finsler-obstruction}
\index{Vector bundle!semipositive curvature obstruction}
Let $F$ carry a holomorphic structure in the neighbourhood of \cref{thm:convex-norm-deformation-obstruction}. If $F^*$ is ample, it admits no strongly pseudoconvex squared Finsler metric with semipositive Kobayashi curvature.
\end{corollary}
\begin{proof}
Apply \cref{cor:ample-semipositive-finsler-smooth-dual} to $E=F^*$. Such a metric would produce a smooth complex norm on $F$ with plurisubharmonic square, contrary to \cref{thm:convex-norm-deformation-obstruction}.
\end{proof}

This conclusion retains the specific Gauduchon degree hypotheses and the nearby holomorphic structure. It does not assert existence of an ample stable family in that neighbourhood. The polar square can be nonsmooth; regularization is an essential separate step.

\subsection{Ample comparison quotients on Hirzebruch surfaces}
\label{subsec:hirzebruch-ample-comparison}
\index{Hirzebruch surface!ample comparison}

To apply \cref{cor:gauduchon-semipositive-finsler-obstruction}, we need ample vector bundles whose duals are close to a split bundle $F_1\oplus F_2$ with $\deg_{\Theta_s}F_s>0$ for suitable test forms $\Theta_s$. The split bundle $F_1^*\oplus F_2^*$ is not ample itself: otherwise its quotient line bundles $F_s^*$ would be ample, hence carry positively curved metrics, and each $F_s$ would have negative degree for every test form. The ample bundles must therefore be obtained by deformation. On a Hirzebruch surface we choose a split bundle $E_0=L_1\oplus L_2$ whose dual summands have positive degree for suitable K\"ahler forms, and we construct an ample bundle $E_{\mathrm{amp}}$ with the same Chern classes as a quotient of a trivial bundle. Both bundles belong to the family of bundles with these Chern classes that are prioritary for the ruling. This family is irreducible, which will allow us to find ample bundles in every neighborhood of $E_0$.

Let $e\geq0$ be an integer, and let
\[
 X=\mathbb F_e=\mathbb P_{\mathbb P^1}(\mathcal O\oplus\mathcal O(e)),
 \qquad \pi\colon X\longrightarrow\mathbb P^1,
\]
be the Hirzebruch surface, where the projective bundle parametrizes lines. Write $S$ for the section defined by the subbundle $\mathcal O(e)$ and $f$ for a fiber of $\pi$, called a ruling fiber; we identify line bundles with divisor classes and write $\mathcal O_X(aS+bf)$. The normal bundle of $S$ is $\mathcal O(-e)$, so
\[
 S^2=-e,\qquad S\cdot f=1,\qquad f^2=0.
\]
The dual of the tautological line bundle is $\mathcal O_X(S)$. A line bundle of degree zero on the ruling fibers is pulled back from $\mathbb P^1$: its direct image has rank one, and the evaluation map from the pullback of the direct image is an isomorphism on each fiber. Consequently $\operatorname{Pic}(X)=\mathbb ZS\oplus\mathbb Zf$.

\Needspace{5\baselineskip}
\begin{lemma}\label{lem:hirzebruch-ruling-sections}
\index{Hirzebruch surface!ruling cohomology}
For integers $a\geq0$ and $b$,
\[
 \pi_*\mathcal O_X(aS+bf)=
 \bigoplus_{i=0}^a\mathcal O_{\mathbb P^1}(b-ie),
 \qquad R^1\pi_*\mathcal O_X(aS+bf)=0.
\]
A class $aS+bf$ is effective exactly when $a,b\geq0$. If $a\geq1$ and $b>ae$, its line bundle is very ample. In particular,
\[
 H=S+(3e+1)f,
 \qquad H^2=5e+2,
 \qquad h^0(X,\mathcal O_X(H))=5e+4.
\]
The canonical class is $K_X=-2S-(e+2)f$.
\end{lemma}
\begin{proof}
Over a trivializing chart of the base, sections of $\mathcal O_X(aS)$ are homogeneous polynomials of degree $a$ in the fiber coordinates. Their transition law gives
\[
 \pi_*\mathcal O_X(aS)=
 \operatorname{Sym}^a(\mathcal O\oplus\mathcal O(-e)),
\]
and twisting by the base gives the formula. The first cohomology of $\mathcal O_{\mathbb P^1}(a)$ on each fiber vanishes, which gives the assertion on the higher direct image. If $a<0$, restriction to a general ruling fiber excludes a nonzero section. For $a\geq0$, the direct-image formula gives a nonzero section exactly when $b\geq0$.

Suppose that $a\geq1$ and $b>ae$. All line bundles in the direct image then have degree at least one. Fiber polynomials of degree $a$ separate points and tangent directions within a fiber. To separate points over different base points, choose a fiber monomial that is nonzero at the first point and a coefficient section vanishing at the second base point but not at the first. At any point, choose such a nonzero monomial and a coefficient section with a simple zero at its base point; the resulting global section vanishes there, has nonzero derivative in a prescribed base direction and has zero vertical derivative. The complete linear system therefore separates points and first jets, so it defines an embedding. The formulas for $H^2$ and $h^0(H)$ follow from the intersection numbers and from the direct image $\mathcal O(3e+1)\oplus\mathcal O(2e+1)$. Finally, the relative Euler sequence gives $K_{X/\mathbb P^1}=-2S-ef$, and adding the pullback of the canonical class of $\mathbb P^1$ proves the last identity.
\end{proof}

Put
\[
 L_1=\mathcal O_X(5S-f),\qquad
 L_2=\mathcal O_X(-S+(12e+5)f),\qquad
 E_0=L_1\oplus L_2.
\]
Their Chern classes satisfy
\[
 c_1(E_0)=4H,\qquad c_2(E_0)=65e+26=13H^2.
\]
An ample bundle with these Chern classes can be constructed as a quotient of a trivial bundle. The main point of the proof is to make the fibers of the associated map to projective space finite, which requires excluding common curves of the defining sections uniformly over the projective parameter.

\Needspace{5\baselineskip}
\begin{proposition}\label{prop:hirzebruch-ample-comparison}
\index{Vector bundle!ample comparison quotient}
There exist linearly independent sections $a_0,\ldots,a_3$ of $\mathcal O_X(H)$ with no common zero. For any such quadruple and a general quadruple $b_0,\ldots,b_3$ of sections of $\mathcal O_X(3H)$, the sequence
\[
 0\longrightarrow\mathcal O_X(-H)\oplus\mathcal O_X(-3H)
 \xrightarrow{(a,b)}\mathcal O_X^{\oplus4}
 \longrightarrow E_{\mathrm{amp}}\longrightarrow0
\]
has a locally free ample cokernel $E_{\mathrm{amp}}$ of rank two, with $c_1(E_{\mathrm{amp}})=4H$ and $c_2(E_{\mathrm{amp}})=13H^2$. Here \emph{general} means outside a proper Zariski-closed subset of the space of quadruples.
\end{proposition}
\begin{proof}
\emph{Step 1: the quadruple $a$.} Since $\mathcal O_X(H)$ is globally generated, a common zero at a fixed point imposes four independent linear conditions on a quadruple $a$. The incidence variety of pairs consisting of a point of the surface and a quadruple vanishing there therefore has dimension at most two less than that of the space of quadruples, and its image in that space is closed by properness. Hence quadruples with a common zero form a proper closed subset. Linear independence is also a nonempty open condition, because $h^0(H)\geq4$, and the two open conditions intersect. Fix such an $a$ and write
\[
 V=H^0(X,\mathcal O_X(3H))^{\oplus4},\qquad N=\dim V.
\]

\emph{Step 2: local freeness.} At a fixed point $x$, the map $(a,b)$ drops rank exactly when $b(x)$ is proportional to the nonzero vector $a(x)$. Since evaluation on $V$ is surjective, this is a condition of codimension three. Its incidence over $X$ has dimension at most $N-1$, and its image in $V$ is closed. Outside that image the cokernel is locally free of rank two.

\emph{Step 3: the map to projective space.} Dualizing the sequence identifies $E_{\mathrm{amp}}^*$ with a subbundle of $(\mathbb C^4)^*\otimes\mathcal O_X$. This inclusion defines a proper morphism
\[
 q_b\colon\mathbb P(E_{\mathrm{amp}}^*)\longrightarrow
 \mathbb P((\mathbb C^4)^*),
 \qquad \mathcal O_{\mathbb P(E_{\mathrm{amp}}^*)}(1)
 =q_b^*\mathcal O_{\mathbb P^3}(1).
\]
For $w=[w_0:\cdots:w_3]$, its fiber is the common zero scheme of
\[
 s_w=\sum_iw_i a_i\in H^0(X,\mathcal O_X(H)),
 \qquad t_w=\sum_iw_i b_i\in H^0(X,\mathcal O_X(3H)).
\]
Linear independence of $a$ gives $s_w\neq0$ for every $w$.

\emph{Step 4: no common curves.} A nonzero effective subdivisor $D$ of $\operatorname{div}(s_w)$ has class $uS+vf$, where
\[
 u\in\{0,1\},\qquad 0\leq v\leq3e+1,
 \qquad (u,v)\neq(0,0).
\]
Write $|D|=\mathbb P H^0(X,\mathcal O_X(D))$ for the complete linear system. For each of these finitely many classes $D_0$, let $T_{D_0}$ parametrize the triples $(w,[\sigma],[\tau])$ with
\[
 [\sigma]\in|D_0|,\qquad [\tau]\in|H-D_0|,
 \qquad [s_w]=[\sigma\tau].
\]
Since products of nonzero sections are nonzero, this is a closed projective parameter scheme. Its fiber over $w$ consists of the finitely many effective subdivisors of the fixed divisor $\operatorname{div}(s_w)$, counted with their possible multiplicities. Thus $\dim T_{D_0}\leq3$.

Vanishing on a nonzero effective divisor $D$ imposes at least four independent conditions on $H^0(X,\mathcal O_X(3H))$. Indeed, choose an integral curve $C$ in its support. Very ampleness gives sections $r_0,r_1$ of $\mathcal O_X(H)$ with nonconstant ratio on $C$. The four cubic monomials in their restrictions are linearly independent, since a linear relation would make the nonconstant ratio algebraic over $\mathbb C$, hence constant.

For fixed $(w,[\sigma],[\tau])$, the map $V\to H^0(X,\mathcal O_X(3H))$, $b\mapsto t_w$, is surjective. The condition that $t_w$ vanish on $D=\operatorname{div}(\sigma)$ therefore has codimension at least four in $V$. This condition is closed in families: the universal multiplication map
\[
 H^0(X,\mathcal O_X(3H-D_0))\otimes\mathcal O_{|D_0|}(-1)
 \longrightarrow H^0(X,\mathcal O_X(3H))\otimes\mathcal O_{|D_0|}
\]
is injective on fibers, so its image is a subbundle, and the condition is the vanishing of the image of $t_w$ in the quotient bundle. Its incidence over $T_{D_0}$ has dimension at most $N-1$, and its image in $V$ is closed.

\emph{Step 5: conclusion.} Choose $b$ outside these finitely many proper closed images and outside the image of Step~2. Then $s_w$ and $t_w$ have no common curve for any $w$, so every fiber of $q_b$ is finite. A proper morphism with finite fibers is finite, and the pullback of $\mathcal O_{\mathbb P^3}(1)$ by a finite morphism is ample; hence $E_{\mathrm{amp}}$ is ample. Finally, the Whitney formula gives
\[
 c(E_{\mathrm{amp}})=\frac1{(1-H)(1-3H)}=1+4H+13H^2.
\]
\end{proof}

\index{Prioritary bundle}
A vector bundle $E$ on $X$ is called \emph{prioritary for the ruling} if $\operatorname{Ext}^2_X(E,E(-f))=0$. Since all ruling fibers are linearly equivalent, this condition does not depend on the chosen fiber, and by Serre duality it is equivalent to
\[
 \operatorname{Hom}_X(E,E(K_X+f))=0.
\]
Prioritary bundles of fixed rank and Chern classes form an irreducible family, by the theorem of Walter recalled in \cref{subsec:hirzebruch-stable-ample-deformations}.

\Needspace{5\baselineskip}
\begin{lemma}\label{lem:hirzebruch-prioritary-comparison}
\index{Prioritary bundle!comparison points}
Both $E_0$ and $E_{\mathrm{amp}}$ are prioritary for the ruling.
\end{lemma}
\begin{proof}
Write $D=K_X+f=-2S-(e+1)f$. The summands of $\operatorname{End}(E_0)(D)$ have classes
\[
 D,\quad D,\quad 4S-(13e+7)f,\quad -8S+(11e+5)f.
\]
Each has a negative coefficient in the basis $S,f$, so by \cref{lem:hirzebruch-ruling-sections} none of them has a nonzero section. This proves the required vanishing of $\operatorname{Hom}$ for $E_0$.

Precomposition with the surjection $\mathcal O_X^4\to E_{\mathrm{amp}}$ gives an injection
\[
 \operatorname{Hom}(E_{\mathrm{amp}},E_{\mathrm{amp}}(D))
 \hookrightarrow H^0(X,E_{\mathrm{amp}}(D))^4.
\]
Twist the defining sequence by $D$. Its middle term has no sections. For $j=1,3$, Serre duality and the formula of \cref{lem:hirzebruch-ruling-sections} give
\[
 \begin{aligned}
 H^1(X,\mathcal O_X(D-jH))&=H^1(X,\mathcal O_X(jH-f))^*,\\
 \pi_*\mathcal O_X(jH-f)&=\bigoplus_{i=0}^j\mathcal O_{\mathbb P^1}(j(3e+1)-1-ie).
 \end{aligned}
\]
The higher direct image vanishes, and every degree in this sum is at least $j(2e+1)-1\geq0$. Hence all these first cohomology groups vanish, and the long exact sequence of the twisted sequence gives $H^0(X,E_{\mathrm{amp}}(D))=0$. This proves the vanishing of $\operatorname{Hom}$ for $E_{\mathrm{amp}}$.
\end{proof}

\Needspace{5\baselineskip}
\begin{proposition}\label{prop:hirzebruch-positive-dual-tests}
\index{Gauduchon degree!Hirzebruch dual tests}
For $F_s=L_s^*$, $s=1,2$, there are K\"ahler forms $\eta_s$ with $\deg_{\eta_s}F_s=1$. Consequently, the split bundle $E_0^*=F_1\oplus F_2$ satisfies the hypotheses of \cref{thm:convex-norm-deformation-obstruction}.
\end{proposition}
\begin{proof}
The integral classes
\[
 A_1=6S+(6e+1)f,\qquad A_2=S+(13e+6)f
\]
are very ample by \cref{lem:hirzebruch-ruling-sections}. Pullbacks of normalized Fubini--Study forms under the corresponding embeddings give K\"ahler forms $\eta_s$ representing $c_1(\mathcal O_X(A_s))$. A direct computation of intersection numbers gives
\[
 c_1(L_1)\cdot A_1=-1,\qquad
 c_1(L_2)\cdot A_2=-1,
\]
so the degrees of the duals are equal to one. Since $n=2$, the closed positive $(1,1)$-forms $\eta_s$ are admissible test forms $\Theta_s$ in \cref{thm:convex-norm-deformation-obstruction}.
\end{proof}

\subsection{Stable extensions approaching the split comparison bundle}
\label{subsec:hirzebruch-stable-extensions}
\index{Hirzebruch surface!stable extensions}

The bundles near $E_0$ that we construct will also be slope-stable for a fixed polarization. The split bundle $E_0$ is unstable, and the ample bundle $E_{\mathrm{amp}}$ need not be stable. Nonsplit extensions of $L_2$ by $L_1$ have the same Chern classes as $E_0$, are stable, and converge to $E_0$ as the extension class tends to zero; they are not ample, however. The next subsection combines them with $E_{\mathrm{amp}}$.

Keep the Hirzebruch surface $X=\mathbb F_e$, its ruling and the line bundles of \cref{subsec:hirzebruch-ample-comparison}. Thus
\[
 L_1=\mathcal O_X(5S-f),\qquad
 L_2=\mathcal O_X(-S+(12e+5)f),\qquad E_0=L_1\oplus L_2.
\]
We use the very ample polarization $P=4H-f=4S+(12e+3)f$. For a torsion-free coherent sheaf $E$ of positive rank on $X$, put
\[
 \deg_PE=c_1(E)\cdot P,
 \qquad \mu_P(E)=\frac{\deg_PE}{\operatorname{rank}E}.
\]
A vector bundle $V$ is \emph{$\mu_P$-stable} if $\mu_P(M)<\mu_P(V)$ for every coherent subsheaf $M\subset V$ with $0<\operatorname{rank}M<\operatorname{rank}V$. Since $\mu_P(L_2)>\mu_P(E_0)$ by the computation below, the split bundle $E_0$ is not $\mu_P$-stable.

\Needspace{5\baselineskip}
\begin{proposition}\label{prop:hirzebruch-stable-extensions}
\index{Vector bundle!slope-stable extensions}
The extension space has dimension
\[
 \dim\operatorname{Ext}^1_X(L_2,L_1)=105e+35.
\]
Every nonsplit extension
\[
 0\longrightarrow L_1\longrightarrow V\xrightarrow{q}L_2
 \longrightarrow0
\]
is a locally free, $\mu_P$-stable bundle with $c_1(V)=4H$ and $c_2(V)=13H^2$.
\end{proposition}
\begin{proof}
Extensions of line bundles are locally free: the quotient is free on a small trivializing neighborhood, so the sequence splits there. Their classes form
\[
 \operatorname{Ext}^1_X(L_2,L_1)
 =H^1(X,\mathcal O_X(6S-(12e+6)f)).
\]
By \cref{lem:hirzebruch-ruling-sections}, the higher direct image of this line bundle vanishes and its direct image is the sum of the seven line bundles $\mathcal O_{\mathbb P^1}(-12e-6-ie)$, $0\leq i\leq6$. Their first cohomology groups have total dimension
\[
 \sum_{i=0}^6(12e+5+ie)=105e+35>0.
\]
The Chern classes follow from the extension sequence. Direct computations of intersection numbers give
\[
 \mu_P(L_1)=40e+11,\qquad \mu_P(V)=40e+14,
 \qquad \mu_P(L_2)=40e+17.
\]

It suffices to test line-bundle subsheaves $M\hookrightarrow V$. Indeed, a rank-one torsion-free subsheaf on a smooth surface has a line bundle as its double dual, and its inclusion extends across the finitely many missing points by Hartogs extension in local bundle frames. The extended map remains injective and has the same first Chern class. This argument does not require the quotient to be locally free.

If the composite $M\to L_2$ is zero, then $M$ maps nontrivially into $L_1$, so $\deg_P M\leq\mu_P(L_1)<\mu_P(V)$. Otherwise its zero divisor $D$ is effective and $M=L_2(-D)$. If $D=0$, the composite is an isomorphism and its inverse splits $q$, contrary to the hypothesis. Hence $D\neq0$, and
\[
 \deg_P M=\mu_P(V)+3-D\cdot P.
\]
Every effective class is $aS+bf$ with $a,b\geq0$, and
\[
 S\cdot P=8e+3,\qquad f\cdot P=4.
\]
For $e\geq1$, every nonzero effective divisor has $P$-degree at least four, which proves the strict slope inequality.

For $e=0$, the only nonzero effective class of $P$-degree at most three is $S$. A lift of multiplication by a nonzero section $s\in H^0(X,\mathcal O_X(S))$ from $L_2(-S)$ to $V$ would make the pullback of the nonzero extension class vanish. That pullback is the image of the class under the map
\[
 H^1(X,\mathcal O_X(6S-6f))
 \xrightarrow{\ \cdot s\ }H^1(X,\mathcal O_X(7S-6f)).
\]
On $X=\mathbb P^1\times\mathbb P^1$, order the factors so that $\mathcal O_X(S)=\mathcal O(1,0)$ and $\mathcal O_X(f)=\mathcal O(0,1)$. The K\"unneth formula identifies this map with
\[
 \bigl[H^0(\mathbb P^1,\mathcal O(6))
 \xrightarrow{\ \cdot s\ }H^0(\mathbb P^1,\mathcal O(7))\bigr]
 \otimes\operatorname{Id}_{H^1(\mathbb P^1,\mathcal O(-6))}.
\]
Multiplication by a nonzero section is injective on global sections, and tensoring with a vector space preserves injectivity. Thus the nonzero extension class cannot pull back to zero. This excludes the possible equality of slopes and proves stability also for $e=0$.
\end{proof}

\Needspace{5\baselineskip}
\begin{proposition}\label{prop:hirzebruch-extension-prioritary}
\index{Prioritary bundle!extensions}
\index{Vector bundle!negative fibre quotient}
Every extension of $L_2$ by $L_1$, including the split extension, is prioritary for the ruling and is not ample. Its restriction to every ruling fiber is $\mathcal O_{\mathbb P^1}(5)\oplus\mathcal O_{\mathbb P^1}(-1)$.
\end{proposition}
\begin{proof}
Put $D=K_X+f=-2S-(e+1)f$. The four line bundles $L_i^*\otimes L_j(D)$ have classes
\[
 D,\quad D,\quad 4S-(13e+7)f,\quad -8S+(11e+5)f.
\]
Each has a negative coefficient in the basis $S,f$, and hence no nonzero section. Applying $\operatorname{Hom}_X(L_i,\cdot)$ to the extension twisted by $D$ gives $\operatorname{Hom}_X(L_i,V(D))=0$ for $i=1,2$. Applying $\operatorname{Hom}_X(\cdot,V(D))$ to the original extension then gives $\operatorname{Hom}_X(V,V(D))=0$. By Serre duality this is the prioritary condition $\operatorname{Ext}^2_X(V,V(-f))=0$.

Restriction to a ruling fiber is exact, because the quotient is locally free. The restricted line bundles have degrees five and minus one, and the restricted extension splits, since
$\operatorname{Ext}^1_{\mathbb P^1}(\mathcal O(-1),\mathcal O(5))
=H^1(\mathbb P^1,\mathcal O(6))=0$.
The quotient of negative degree on this curve excludes ampleness.
\end{proof}

\Needspace{5\baselineskip}
\begin{proposition}\label{prop:hirzebruch-explicit-stable-degeneration}
\index{Dolbeault operator!smooth convergence}
\index{Vector bundle!stable degeneration}
On the smooth bundle underlying $E_0$ there is a holomorphic one-parameter family of Dolbeault operators $\bar\partial_t=\bar\partial_0+tA$, $t\in\mathbb C$, such that for $t\ne0$ the holomorphic bundle defined by $\bar\partial_t$ is $\mu_P$-stable and prioritary, but not ample. As $t\to0$, these operators converge to $\bar\partial_0$ in the $C^m$ norm for every integer $m\geq0$.
\end{proposition}
\begin{proof}
Choose a nonzero extension class and a smooth Dolbeault representative $\eta\in\Omega^{0,1}(X,L_2^*\otimes L_1)$ of it, so that $\bar\partial\eta=0$. On the smooth direct sum, set
\[
 A=\begin{pmatrix}0&\eta\\0&0\end{pmatrix}.
\]
The covariant Dolbeault derivative of $A$ is zero and $A\wedge A=0$, so $\bar\partial_t^2=0$. Adding the Dolbeault operator of the base $\mathbb C$ of the parameter $t$ also gives an integrable operator on $X\times\mathbb C$, because the family depends holomorphically on $t$. This defines a holomorphic family of bundles whose extension class is $t[\eta]$. It is nonsplit for $t\neq0$, so the two preceding propositions apply. For a fixed smooth background connection and fixed metrics,
\[
 \|\bar\partial_t-\bar\partial_0\|_{C^m}
 =|t|\,\|A\|_{C^m}\longrightarrow0,
\]
and the norm $\|A\|_{C^m}$ is finite for every $m$ because $X$ is compact.
\end{proof}

% Stable-extension attribution: Liu--Wan 2609.37309v1, Proposition 4.4, pp. 20--21.
Thus stability and ampleness have been realized separately in the prioritary family of $E_0$: the nonsplit extensions are stable but not ample, while $E_{\mathrm{amp}}$ is ample. The deformation argument of \cref{thm:hirzebruch-stable-ample-approximation} joins the stable locus to the ample one and provides stable ample bundles arbitrarily close to $E_0$.

\subsection{Deforming to stable ample bundles}
\label{subsec:hirzebruch-stable-ample-deformations}
\index{Hirzebruch surface!stable ample approximation}

Use the Hirzebruch surface and the data $H,P,L_1,L_2,E_0$ of
\cref{subsec:hirzebruch-ample-comparison,subsec:hirzebruch-stable-extensions}.
The ample quotient and a nonsplit extension give two prioritary bundles
with the same Chern classes. The former is ample and the latter is stable.
To obtain bundles with both properties near $E_0$, we first prove the
openness assertions needed in a parameter family.

Let $T$ be a smooth complex algebraic variety. A flat family below means
a coherent sheaf on $X\times T$ that is flat over $T$.

\Needspace{6\baselineskip}
\begin{lemma}\label{lem:hirzebruch-family-local-freeness}
\index{Vector bundle!local freeness in families}
If a flat family $\mathcal E$ has a locally free fibre at $t_0\in T$,
then $\mathcal E$ is a vector bundle on $X\times T'$ for some Zariski-open
neighbourhood $T'$ of $t_0$.
\end{lemma}
\begin{proof}
At a point of $X\times\{t_0\}$, lift a local basis of the fibre to
$\mathcal E$. Nakayama's lemma makes the resulting map from a finite
free module surjective. Tensor its kernel sequence with the residue field
of $t_0$. Flatness over $T$ makes the first Tor term zero, so the kernel
restricts to the kernel of the chosen fibre isomorphism, namely zero.
Nakayama's lemma kills the kernel at that point. Thus $\mathcal E$ is
locally free near every point of the central fibre. Its non-locally-free
locus is closed. The projection $X\times T\to T$ is proper, so the image
of that locus is closed and misses $t_0$. Remove this image.
\end{proof}

\Needspace{6\baselineskip}
\begin{lemma}\label{lem:hirzebruch-family-ample-open}
\index{Vector bundle!openness of ampleness}
For a vector bundle $\mathcal E$ on $X\times T$, the parameters $t$
for which $\mathcal E_t$ is ample form a Zariski-open subset of $T$.
\end{lemma}
\begin{proof}
Put $Y=\mathbb P_{X\times T}(\mathcal E^*)$, with the line convention,
and let $p:Y\to T$. This morphism is projective and flat. Write
$\mathcal L=\mathcal O_Y(1)$. If $\mathcal E_{t_0}$ is ample,
$\mathcal L_{t_0}$ is ample. Choose $k$ so that its $k$th power is very
ample and has zero higher cohomology. Serre vanishing, semicontinuity and
cohomology and base change show, after shrinking around $t_0$, that
$\mathcal V=p_*\mathcal L^k$ is locally free and commutes with fibre
restriction. The evaluation $p^*\mathcal V\to\mathcal L^k$ is surjective:
its cokernel vanishes on the central fibre, and its support has closed
image under the proper map $p$.

Evaluation gives a proper morphism
\[
 g:Y\longrightarrow\mathbb P_T(\mathcal V^*),
 \qquad g^*\mathcal O(1)=\mathcal L^k,
\]
whose central fibre is a closed immersion. The quasi-finite locus of $g$
is open. Its complement has closed image in $T$ by properness and misses
$t_0$, so another shrinking makes $g$ proper and quasi-finite, hence
finite. For a finite morphism, pushforward commutes with base change.
The cokernel of
\[
 \mathcal O_{\mathbb P_T(\mathcal V^*)}\longrightarrow g_*\mathcal O_Y
\]
therefore vanishes on the central fibre. Its support projects to a closed
subset of $T$ missing $t_0$. After removing that subset, the displayed map
is surjective, so $g$ is a closed immersion. Consequently $\mathcal L^k$
is very ample on every fibre, and $\mathcal E_t$ is ample there.
\end{proof}

\Needspace{6\baselineskip}
\begin{lemma}\label{lem:hirzebruch-family-stability-open}
\index{Vector bundle!openness of slope-stability}
For a rank-two vector bundle family $\mathcal E$ with fixed first Chern
class, the $\mu_P$-stable fibre locus is Zariski open in $T$.
\end{lemma}
\begin{proof}
Fix a stable fibre at $t_0$ and set
$\nu=c_1(\mathcal E_t)\cdot P/2$, which is constant. Choose $k\geq1$
so that $\mathcal E_{t_0}^*(kP)$ is globally generated and has zero
higher cohomology. The same Serre, base-change and evaluation argument
gives, on a neighbourhood of $t_0$, a fibrewise injection
\[
 \mathcal E_t\longrightarrow\mathcal O_X(kP)^{\oplus N}.
\]
By the rank-one double-dual argument in the proof of
\cref{prop:hirzebruch-stable-extensions}, a destabilizing subsheaf yields
a nonzero map from a line bundle $M=\mathcal O_X(aS+bf)$ with
$\deg_P M\geq\nu$. At least one component of its composite with the
displayed injection is nonzero. The effective-cone calculation gives
\[
 a\leq4k,\qquad b\leq(12e+3)k.
\]
Since $\deg_P M=(8e+3)a+4b$, the slope inequality also gives
\[
 a\geq\frac{\nu-4(12e+3)k}{8e+3},
 \qquad b\geq\frac{\nu-4k(8e+3)}4.
\]
Only finitely many integral pairs satisfy these four bounds.
For each pair with $\deg_P M\geq\nu$, semicontinuity makes
\[
 \{t:H^0(X,\mathcal E_t\otimes M^*)\neq0\}
\]
closed. Each such subset misses $t_0$. Conversely a nonzero map from
such an $M$ is injective as a sheaf map and violates slope-stability.
Thus the unstable locus on this neighbourhood is exactly the finite
union of these closed subsets. Removing the union proves the claim.
\end{proof}

\Needspace{6\baselineskip}
\begin{lemma}\label{lem:dense-family-smooth-dolbeault-convergence}
\index{Dolbeault operator!algebraic family convergence}
\index{Vector bundle!parallel transport in families}
Let $T$ be smooth and irreducible, let $U\subset T$ be a nonempty
Zariski-open subset, and let $\mathcal E$ be a vector bundle on
$X\times T$. For any $t_0\in T$, there are $t_j\in U$ tending to $t_0$
in the analytic topology and smooth complex-linear identifications with
$\mathcal E_{t_0}$ under which the fibre Dolbeault operators converge
in the $C^m$ norm for every integer $m\geq0$. The dual operators converge
in the same sense.
\end{lemma}
\begin{proof}
The proper algebraic closed subset $T\setminus U$ has empty interior
in the complex-analytic topology. In an analytic coordinate ball $B$
around $t_0$, choose $t_j\in U\cap B$ with $t_j\to t_0$.
Choose a smooth complex connection on $\mathcal E|_{X\times B}$.
Parallel transport along radial paths in $B$, keeping $x\in X$ fixed,
gives identifications
\[
 \Phi_t:\mathcal E_{t_0}^{\mathrm{sm}}\longrightarrow
          \mathcal E_t^{\mathrm{sm}},\qquad \Phi_{t_0}=\operatorname{id},
\]
depending smoothly on $(x,t)$. The superscript $\mathrm{sm}$ denotes
the underlying smooth complex bundle. These identifications need not
be holomorphic.
Transport the Dolbeault operators by these maps and write
\[
 \bar\partial_t=\Phi_t^{-1}\bar\partial_{\mathcal E_t}\Phi_t,
 \qquad A_t=\bar\partial_t-\bar\partial_{\mathcal E_{t_0}}.
\]
Their common Leibniz rule makes $A_t$ an endomorphism-valued $(0,1)$-form.
Its coefficients are smooth in $(x,t)$ and $A_{t_0}=0$. Fix background
metrics and a smooth connection on the central bundle. Compactness of
$X$ and continuity of every finite collection of derivatives give
\[
 \|A_{t_j}\|_{C^m(X)}\longrightarrow0
 \qquad\text{for every integer }m\geq0.
\]
The dual identifications are $(\Phi_t^{-1})^*$. In transported dual
frames the perturbation is $-A_t^{\mathsf T}$, so the same convergence
holds for the dual Dolbeault operators.
\end{proof}

\index{Prioritary bundle!stack irreducibility}
We use Walter's prioritary-stack theorem: on a birationally ruled
surface, a nonempty stack of prioritary torsion-free sheaves with fixed
rank $r\geq2$ and numerical Chern classes is smooth and irreducible.
Here prioritary means $\operatorname{Ext}^2(E,E(-f))=0$ for every ruling
fibre. On $X$ all ruling fibres are linearly equivalent.
This theorem concerns the full prioritary stack and has no semistability
hypothesis.

\Needspace{6\baselineskip}
\begin{theorem}\label{thm:hirzebruch-stable-ample-approximation}
\index{Vector bundle!stable ample approximation}
For every $e\geq0$, there are integrable Dolbeault operators on the fixed
smooth bundle underlying $E_0$ that converge to its split operator in
every finite $C^m$ norm and define ample $\mu_P$-stable bundles of rank
two with $c_1=4H$ and $c_2=13H^2$. The induced dual operators converge
to the split dual operator in every finite $C^m$ norm.
\end{theorem}
\begin{proof}
Let $\mathcal P$ be the prioritary stack with invariants
$(2,4H,13H^2)$. It contains $E_0$ by
\cref{prop:hirzebruch-extension-prioritary}. Walter's theorem makes it
smooth and irreducible.
The loci $\mathcal A$ of ample vector bundles and $\mathcal S$ of
$\mu_P$-stable vector bundles are open by the three family lemmas.
These lemmas apply on smooth local charts of $\mathcal P$; the first
allows restriction to locally free families near a vector-bundle point.
The ample bundle of \cref{prop:hirzebruch-ample-comparison} is prioritary
by \cref{lem:hirzebruch-prioritary-comparison}, so $\mathcal A$ is nonempty.
A nonsplit extension is stable and prioritary by
\cref{prop:hirzebruch-stable-extensions,prop:hirzebruch-extension-prioritary},
so $\mathcal S$ is nonempty. Irreducibility makes
$\mathcal U=\mathcal A\cap\mathcal S$ nonempty and dense.

Choose a smooth local chart
$q:T\to\mathcal P$ of finite type over $\mathbb C$, with a point $t_0$
representing $E_0$ and its flat family $\mathcal E$ on $X\times T$.
Fix an identification $\mathcal E_{t_0}\simeq E_0$.
The variety $T$ is smooth because the stack and chart are smooth.
Shrink to an irreducible neighbourhood of $t_0$ and use
\cref{lem:hirzebruch-family-local-freeness} to make $\mathcal E$ a vector
bundle on the whole product. The inverse image $q^{-1}(\mathcal U)$ is
dense: any nonempty open $T'\subset T$ has nonempty open image $q(T')$
under the smooth map, and that image meets $\mathcal U$.
Apply \cref{lem:dense-family-smooth-dolbeault-convergence} to this dense
open set. Each transported operator defines a bundle isomorphic to the
corresponding ample stable fibre. The fixed stack invariants give its
Chern classes, and the lemma also gives convergence of the dual operators.
\end{proof}

\Needspace{6\baselineskip}
\begin{corollary}\label{cor:hirzebruch-stable-ample-finsler-obstruction}
\index{Kobayashi curvature!stable ample obstruction}
\index{Finsler metric!stable ample obstruction}
For every $e\geq0$, there is an ample $\mu_P$-stable rank-two bundle $E$
on $X$ with $c_1(E)=4H$ and $c_2(E)=13H^2$ that admits no smooth strongly
pseudoconvex squared Finsler metric with semipositive Kobayashi curvature.
\end{corollary}
\begin{proof}
The two dual lines of $E_0^*$ have the positive individual Gauduchon
tests of \cref{prop:hirzebruch-positive-dual-tests}. The preceding theorem
gives an ample stable bundle $E$ whose dual Dolbeault operator is as
close to that split operator in $C^4$ as required. Apply
\cref{cor:gauduchon-semipositive-finsler-obstruction} to $F=E^*$.
\end{proof}

% Exact inputs: Walter alg-geom/9311005v1, Proposition 2, p. 2.
% Approximation: Liu--Wan 2609.37309v1, Theorem 4.1, p. 17; proof, p. 22.
% Curvature conclusion: Liu--Wan v1, Theorem 1.1, p. 2; proof, p. 23.

\subsection{A fixed Stein hypersurface for all Finsler metrics}
\label{subsec:fixed-stein-finsler-hypersurface}
\index{Stein manifold!fixed Finsler test hypersurface}

\Cref{cor:hirzebruch-stable-ample-finsler-obstruction} shows that no strongly pseudoconvex Finsler metric on $E$ has semipositive Kobayashi curvature, but it does not locate the failure. For a single metric, a point where the curvature is not semipositive gives a local hypersurface through that point on which the metric has a nondegenerate critical point with many positive Levi directions. Holomorphic approximation on a Stein manifold allows one hypersurface to serve for all metrics at once.

In this subsection, $X$ is a connected smooth projective variety of dimension $n\geq1$, and $E\to X$ is an algebraic vector bundle of rank $r\geq2$, viewed as a holomorphic bundle. We use the squared Finsler metrics and the convention for Kobayashi curvature of \cref{subsec:polar-finsler-convexity}. A compact subset $K$ of a Stein manifold $Y$ is \emph{holomorphically convex} if it coincides with its holomorphic hull $\{y\in Y:|f(y)|\leq\sup_K|f|\text{ for all }f\in\mathcal O(Y)\}$. We use the classical Oka--Weil theorem in the following form: a function holomorphic near a holomorphically convex compact subset of a Stein manifold can be approximated uniformly there by global holomorphic functions. We also use polynomial Runge approximation in one variable on a union of two disjoint closed disks.

\Needspace{6\baselineskip}
\begin{lemma}\label{lem:separated-holomorphic-convex-compacts}
\index{Holomorphic convexity!separated compact sets}
Let $K,L$ be holomorphically convex compact subsets of a Stein manifold $Y$. If a global holomorphic function maps them into disjoint closed disks in $\mathbb C$, then $K\cup L$ is holomorphically convex.
\end{lemma}
\begin{proof}
Write $h$ for the function and $\Delta_K,\Delta_L$ for the disks. Their union is polynomially convex, so testing with polynomials in $h$ shows that $h$ maps the holomorphic hull of $K\cup L$ into $\Delta_K\cup\Delta_L$. Suppose that $p\notin K\cup L$ and $h(p)\in\Delta_K$. Holomorphic convexity of $K$ gives $a\in\mathcal O(Y)$ with $|a(p)|>\sup_K|a|$. Choose polynomials tending uniformly to one on $\Delta_K$ and to zero on $\Delta_L$. Their compositions with $h$, multiplied by $a$, converge to $a$ at $p$ and on $K$, and to zero on $L$. A sufficiently good approximation therefore separates $p$ from $K\cup L$. The case $h(p)\in\Delta_L$ is identical, with $K$ and $L$ interchanged. Thus no point outside $K\cup L$ belongs to its hull.
\end{proof}

\Needspace{6\baselineskip}
\begin{lemma}\label{lem:dilated-countable-holomorphic-tests}
\index{Holomorphic approximation!countable dilated tests}
Let $U$ be a connected smooth affine variety and let $q\geq0$ be an integer. Put $B=U\times\mathbb C^*\times\mathbb C^q$, with coordinates $(x,t,s)$, and for $\lambda\in\mathbb C^*$ let $d_\lambda(x,t,s)=(x,\lambda t,\lambda s)$. For each $j\geq1$, let $C_j\subset B$ be a holomorphically convex compact set, let $g_j$ be holomorphic near $C_j$, and let $\delta_j>0$. Suppose that $t(C_j)$ lies in the disk $|t-t_j|\leq|t_j|/4$ for some $t_j\neq0$. Then there are $\lambda_j\in\mathbb C^*$ and $f\in\mathcal O(B)$ such that
\[
 \sup_{d_{\lambda_j}(C_j)}
 |f-\lambda_j g_j\circ d_{\lambda_j}^{-1}|<|\lambda_j|\delta_j
 \qquad(j\geq1).
\]
\end{lemma}
\begin{proof}
Choose an exhaustion $(L_j)_{j\geq0}$ of the affine, hence Stein, manifold $B$ by holomorphically convex compact sets with $L_j\subset\operatorname{int}L_{j+1}$. We construct $f_j\in\mathcal O(B)$ and compact sets $K_j$ inductively, starting with $f_0=0$ and $K_0=L_0$.

At stage $j\geq1$, choose $M>0$ with $t(K_{j-1})\subset\{|t|\leq M\}$, and choose $R>4M/3$. Dilation by $\lambda_j=R/t_j$ moves the $t$-projection of $C_j$ into the disk $|t-R|\leq R/4$, which is disjoint from $|t|\leq M$. For the rest of the induction we replace $C_j$, $g_j$ and $\delta_j$ by the dilated compact set $d_{\lambda_j}(C_j)$, the function $\lambda_jg_j\circ d_{\lambda_j}^{-1}$ and the tolerance $|\lambda_j|\delta_j$, keeping the same notation. \Cref{lem:separated-holomorphic-convex-compacts}, applied with the function $t$, makes $K_{j-1}\cup C_j$ holomorphically convex.

On disjoint neighborhoods of $K_{j-1}$ and $C_j$, prescribe $f_{j-1}$ and $g_j$, respectively. Oka--Weil approximation gives $f_j\in\mathcal O(B)$ with
\[
 \sup_{K_{j-1}}|f_j-f_{j-1}|<\epsilon_j,
 \qquad \sup_{C_j}|f_j-g_j|<\epsilon_j,
 \qquad
 \epsilon_j=2^{-j-2}\min(1,\delta_1,\ldots,\delta_j).
\]
Let $K_j$ be a member of the fixed exhaustion containing $K_{j-1}\cup C_j\cup L_j$ in its interior. The interiors of the sets $K_j$ exhaust $B$, and the summable errors make $f_j$ converge locally uniformly to some $f\in\mathcal O(B)$. For a fixed $j$, every later approximation controls its error on $C_j$, since $C_j\subset K_{m-1}$ for $m>j$. Consequently
\[
 \sup_{C_j}|f-g_j|
 \leq\sum_{m\geq j}\epsilon_m
 \leq2^{-j-1}\delta_j<\delta_j.
\]
Returning to the original notation gives the stated tolerances.
\end{proof}

\Needspace{6\baselineskip}
\begin{theorem}\label{thm:fixed-stein-finsler-hypersurface}
\index{Finsler metric!fixed hypersurface test}
\index{Kobayashi curvature!critical hyperplane test}
\index{Critical point!holomorphic quadratic jet}
\index{Levi index!fixed hypersurface test}
Suppose that $E$ admits no smooth strongly pseudoconvex squared Finsler metric with semipositive Kobayashi curvature. Then there is a connected smooth Stein hypersurface $\mathcal S\subset E^\times$, locally closed in the analytic topology, with the following property: for every smooth strongly pseudoconvex squared Finsler metric $G$ on $E$, the restriction $G|_{\mathcal S}$ has a nondegenerate critical point with at least $r$ positive Levi eigenvalues, and hence with at most $n-1$ negative ones.
\end{theorem}
\begin{proof}
\emph{Step 1: coordinates.} An algebraic frame trivializes $E$ over a nonempty affine open subset $U\subset X$. Since $X$ is smooth and connected, it is irreducible, so $U$ is dense and connected. Write the fiber coordinates as $(w,t,s)\in\mathbb C\times\mathbb C\times\mathbb C^{r-2}$, and put $B=U\times\mathbb C^*\times\mathbb C^{r-2}$. We seek $\mathcal S$ as a graph $w=f(x,t,s)$ over $B$; such a graph avoids the zero section because $t\ne0$.

\emph{Step 2: a local test for one metric.} Fix a candidate metric $G$. Since its Kobayashi curvature is not semipositive, its horizontal Schur complement has a positive direction somewhere. This condition is open, so it holds at a point $v_0=(x_0,w_0,t_0,s_0)$ with $x_0\in U$, $t_0\neq0$ and $G_w(v_0)\neq0$. Indeed, the first two conditions define a dense open set, and the zero set of $G_w$ has empty interior, since vanishing of $G_w$ on an open set would give $G_{w\bar w}=0$ there, contradicting the positivity of the vertical Hessian.

Let $H$ be the Levi-orthogonal complement of the vertical tangent space $V$ at $v_0$. The Euler identities in the proof of \cref{lem:polar-finsler-osculating-section} give
\[
 K:=\ker\partial G(v_0)=H\oplus(V\cap K),
 \qquad \dim_{\mathbb C}(V\cap K)=r-1.
\]
Indeed, $\partial G$ vanishes on $H$, and its value on the vertical radial vector is $G(v_0)>0$. The two summands are Levi orthogonal, the vertical summand is positive definite, and the horizontal summand has a positive direction. Thus the Levi form of $G$ restricted to $K$ has at least $r$ positive eigenvalues.

Since $G_w(v_0)\neq0$, the space $K$ is the graph of a linear map defined on $T_{b_0}^{1,0}B$, where $b_0=(x_0,t_0,s_0)$. Choose a local holomorphic function $g$ near $b_0$ with $g(b_0)=w_0$ whose graph has tangent space $K$ at $v_0$. Then $F=G\circ\operatorname{Graph}(g)$ has a critical point at $b_0$, and its Levi form there is the restriction of the Levi form of $G$ to $K$. Keep this first jet of $g$ fixed and vary its quadratic jet. The mixed second derivatives of $F$ do not change, while in local coordinates $z$ on $B$ the pure second derivatives satisfy
\[
 F_{\alpha\beta}(b_0)=A_{\alpha\beta}+G_w(v_0)g_{\alpha\beta}(b_0),
\]
where $A$ is fixed. Since $G_w(v_0)\ne0$, we can prescribe the pure quadratic part of $F$ to be $\tau\operatorname{Re}\sum_\alpha z_\alpha^2$, whose real Hessian has eigenvalues $\pm2\tau$. If $\tau$ is larger than half the operator norm of the real Hessian of the fixed mixed part, the full real Hessian is invertible by the Neumann-series criterion. The critical point is then nondegenerate, with the same Levi form.

\emph{Step 3: stability of the test.} Choose compact neighborhoods $C'\subset\operatorname{int}C$ of $b_0$ on which $g$ is defined, with $C$ holomorphically convex and $t(C)\subset\{|t-t_0|\leq|t_0|/4\}$. For a closed affine embedding of $B$, the intersection of $B$ with a small closed Euclidean ball centered at $b_0$ is such a set $C$, since a closed Euclidean ball is convex and hence polynomially convex. Fix a smooth Hermitian background metric on $E$ and its compact unit sphere bundle $\Sigma$. By homogeneity, the $C^2$ distance between metrics on $\Sigma$ controls their $C^2$ distance on each fixed compact subset of $E^\times$. Persistence of nondegenerate critical points and openness of the conditions on the Levi eigenvalues therefore give a neighborhood $\mathcal V_G$ of $G|_\Sigma$ for the $C^2$ topology and a number $\delta_G>0$ with the following property. For every metric whose restriction to $\Sigma$ lies in $\mathcal V_G$ and every holomorphic function within $\delta_G$ of $g$ uniformly on $C$, the restriction of the metric to the graph of the function still has a nondegenerate critical point over $C'$ with at least $r$ positive Levi eigenvalues. Here the Cauchy estimates on $C'$ turn uniform closeness on $C$ into $C^2$ closeness on $C'$.

\emph{Step 4: one hypersurface for all metrics.} The restrictions of the candidate metrics to $\Sigma$ form a subspace of the separable metrizable space $C^2(\Sigma)$. This subspace is second countable, hence Lindel\"of, so the cover by the neighborhoods $\mathcal V_G$ has a countable subcover $(\mathcal V_j)_{j\geq1}$, with associated data $C_j,g_j,\delta_j$. For $D_\lambda(x,w,t,s)=(x,\lambda w,\lambda t,\lambda s)$,
\[
 G\circ D_\lambda=|\lambda|^2G.
\]
Dilation therefore preserves critical points of restrictions to graphs, their real nondegeneracy and their Levi inertia, for every candidate metric. It maps the graph of a function $g$ over $C$ to the graph of $\lambda g\circ d_\lambda^{-1}$ over $d_\lambda(C)$, with $d_\lambda$ as in \cref{lem:dilated-countable-holomorphic-tests}. A function within $|\lambda|\delta_j$ of $\lambda g_j\circ d_\lambda^{-1}$ on $d_\lambda(C_j)$ pulls back under the dilation to a function within $\delta_j$ of $g_j$ on $C_j$. In particular, the metric neighborhood $\mathcal V_j$ needs no change.

Applying \cref{lem:dilated-countable-holomorphic-tests}, we obtain one function $f\in\mathcal O(B)$ that satisfies all the dilated tests. Its graph $\mathcal S$ is biholomorphic to the connected affine Stein manifold $B$, and it is smooth and closed in the open subset $\{x\in U,\ t\neq0\}$ of $E^\times$. Every candidate metric belongs to some $\mathcal V_j$, so its restriction to $\mathcal S$ has the required nondegenerate critical point. Since $\dim_{\mathbb C}\mathcal S=n+r-1$, at least $r$ positive Levi eigenvalues leave at most $n-1$ negative ones.
\end{proof}

\Needspace{10\baselineskip}
\begin{corollary}\label{cor:hirzebruch-fixed-stein-hypersurface}
\index{Hirzebruch surface!fixed Stein hypersurface}
For each $e\geq0$, the bundle $E$ of \cref{cor:hirzebruch-stable-ample-finsler-obstruction} can be chosen so that $E^\times$ contains a connected smooth Stein hypersurface on which the restriction of every smooth strongly pseudoconvex squared Finsler metric on $E$ has a nondegenerate critical point with at least two positive Levi eigenvalues and at most one negative Levi eigenvalue.
\end{corollary}
\begin{proof}
The bundle constructed in the proof of \cref{thm:hirzebruch-stable-ample-approximation} is a fiber of an algebraic family on a chart of the prioritary stack, so it is algebraic. Apply \cref{thm:fixed-stein-finsler-hypersurface} with $n=r=2$, using the curvature obstruction of \cref{cor:hirzebruch-stable-ample-finsler-obstruction}.
\end{proof}

% Source reconstruction: Liu--Wan 2609.37309v1, Proposition 5.1, pp. 24--27.
% Algebraicity is explicit here; the general holomorphic-to-algebraic GAGA step is not needed.
% Classical accepted inputs: Oka--Weil, one-variable Runge, affine Stein theory.
% The complement-convexity application and Fulton primary proof are not certified here.

\subsection{Concave neighborhoods and dual norms}
\label{subsec:concave-neighbourhood-dual-norms}
\index{Concave neighborhoods!dual support norms}
\index{Convex norm!supporting balls}
\index{Plurisubharmonic function!dual support norms}

A relatively compact neighborhood $\Omega$ of the zero section of a vector bundle $E$ determines a norm on $E^*$, the support function $\xi\mapsto\max_{v\in\overline\Omega\cap E_x}|\xi(v)|$ of its fibers. If the boundary of $\Omega$ is sufficiently concave in the horizontal directions, holomorphic sections touching the boundary from inside show that this support norm has plurisubharmonic square, and supporting balls give it uniform convexity in the fibers. Consequently, when the dual norms are obstructed as in \cref{thm:convex-norm-deformation-obstruction}, such neighborhoods cannot exist. In terms of exhaustion functions, this bounds the convexity of the complement of the zero section in the projective completion of $E$.

In this subsection, $E\to X$ is a holomorphic vector bundle of rank $r\geq1$, and $n\geq1$. Fix a smooth Hermitian metric $k$ on $E$, with real inner product $k_{\mathbb R}=\operatorname{Re}k$ and dual norm $|\cdot|_{k^*}$. For a smooth real function $f$, write $\mathcal L_f$ for its complex Hessian. Let $\Omega$ be an open set with smooth boundary, defined near $\partial\Omega$ by $f<0$ with $\mathrm df\ne0$ on $\partial\Omega$. We say that $\partial\Omega$ has \emph{at least $n$ negative Levi directions} if, at every boundary point, $\mathcal L_f$ has at least $n$ negative eigenvalues on $\ker\partial f$.

\Needspace{7\baselineskip}
\begin{lemma}\label{lem:ball-support-uniform-convexity}
\index{Convex norm!uniform strong convexity}
Fix $\delta,B>0$. On a finite-dimensional Hermitian vector space, put
\[
 s_c(\xi)=|\xi(c)|+\delta|\xi|
\]
for covectors $\xi$ and vectors $c$ with $|c|\leq B$. There is $\gamma>0$, depending only on $\delta$ and $B$, such that $s_c^2-\gamma|\xi|^2$ is real convex for every such $c$.
\end{lemma}
\begin{proof}
For $0<\tau\leq1$, replace $s_c$ by the norm
\[
 s_{c,\tau}(\xi)=\sqrt{|\xi(c)|^2+\tau|\xi|^2}+\delta|\xi|,
\]
which is smooth off zero. Its Lipschitz constant is at most $L=\sqrt{B^2+1}+\delta$, uniformly in $c$ and $\tau$. All derivatives below are real derivatives on the dual space. At $\xi\ne0$, set $u=\xi/|\xi|$ and decompose a vector as $w=\alpha u+b$ with $b$ orthogonal to $u$ for the real inner product. Write $s=s_{c,\tau}$ and $A=\mathrm ds_\xi(w)$. Since the first summand of $s$ is convex and $s\geq\delta|\xi|$,
\[
 D^2s(w,w)\geq\delta\frac{|b|^2}{|\xi|},\qquad
 D^2(s^2)(w,w)\geq2A^2+2\delta^2|b|^2.
\]
Homogeneity gives $\mathrm ds_\xi(u)=s(\xi)/|\xi|\geq\delta$. Since $|\mathrm ds_\xi(b)|\leq L|b|$, we have
\[
 |w|^2\leq\frac{2A^2}{\delta^2}
       +\left(1+\frac{2L^2}{\delta^2}\right)|b|^2.
\]
Take $\gamma=\delta^4/(\delta^2+2L^2)$. Then $\gamma\leq\delta^2/2$ and $\gamma(1+2L^2/\delta^2)=\delta^2$, so $D^2(s^2)(w,w)\geq2\gamma|w|^2$.

The homogeneous function $s^2-\gamma|\xi|^2$ of degree two extends as a $C^1$ function at zero with differential zero. Its restriction to every real line is convex, including at zero, because its second derivative is nonnegative away from zero and its first derivatives match there. Thus it is convex on the whole space. Finally,
\[
 0\leq s_{c,\tau}(\xi)-s_c(\xi)\leq\sqrt\tau\,|\xi|,
\]
and the locally uniform limit as $\tau\downarrow0$ proves the assertion with the same $\gamma$.
\end{proof}

\Needspace{10\baselineskip}
\begin{theorem}\label{thm:concave-neighbourhood-dual-norm}
\index{Concave neighborhoods!smooth dual norms}
Suppose that a relatively compact neighborhood $\Omega\Subset E$ of the zero section has smooth boundary with at least $n$ negative Levi directions. Then
\[
 p(x,\xi)=\max_{v\in\overline\Omega\cap E_x}|\xi(v)|
\]
is a continuous complex norm on $E^*$ whose square $Q=p^2$ is plurisubharmonic on the total space of $E^*$. For some $\gamma>0$, each function $Q(x,\cdot)-\gamma|\cdot|_{k^*}^2$ is real convex. If $E$ is ample, $E^*$ also admits a smooth complex norm with strictly plurisubharmonic square off the zero section.
\end{theorem}
\begin{proof}
\emph{Step 1: basic bounds.} Put $A_x=\overline\Omega\cap E_x$. By compactness of $X$ there are $0<a_0<b_0$ such that, in every fiber, the closed $k$-ball of radius $a_0$ is contained in $\Omega$ and $\overline\Omega$ lies in the open ball of radius $b_0$. The maximum defining $p$ is attained, and
\[
 a_0|\xi|_{k^*}\leq p(x,\xi)\leq b_0|\xi|_{k^*}.
\]
As a supremum of absolute values of linear functionals, $p$ is complex homogeneous and satisfies the triangle inequality.

\emph{Step 2: support points.} Fix $\xi_0\ne0$ and a maximizing point $v_0\in A_{x_0}$. It lies on $\partial\Omega$, since at an interior point a small vertical displacement would increase $|\xi_0(v)|$. Multiply $\xi_0$ by a complex number of modulus one to obtain $\eta$ with $\eta(v_0)=p(x_0,\xi_0)=a>0$. Then $\ell=\operatorname{Re}\eta$ satisfies $\ell\leq a$ on $A_{x_0}$. Let $f$ be a defining function as in the definition of negative Levi directions, so that $\Omega=\{f<0\}$ near $\partial\Omega$ and $\mathrm df\ne0$ on $\partial\Omega$.

Let $V=T^{1,0}_{v_0}E_{x_0}$ be the vertical tangent space and $T=\ker\partial f(v_0)$. Choose an $n$-dimensional complex subspace $W\subset T$ on which $\mathcal L_f$ is negative definite. We first show that the vertical differential $\mathrm d_vf(v_0)$ is nonzero. Otherwise $V\subset T$, and counting dimensions gives $W\cap V\ne\{0\}$. A vector in this intersection has negative vertical Levi form, and the two corresponding real directions show that the real vertical Hessian of $f$ has a negative direction. Its negative cone is open, hence contains a direction $w$ with $\ell(w)>0$. Since the vertical first derivative vanishes, Taylor expansion gives $f(v_0+tw)<0$ for small $t>0$, contradicting $\ell\leq a$ on $A_{x_0}$.

The boundary of $A_{x_0}$ is therefore smooth at $v_0$. Since the supporting functional $\ell$ attains its maximum on $A_{x_0}$ there,
\[
 \mathrm d_vf(v_0)=c\,\mathrm d\ell,\qquad c>0;
\]
the sign follows because an inward direction along which $\ell$ increases would contradict the maximality. On the affine hyperplane $\{\ell=a\}$ of the fiber, the restriction of $f$ has a local minimum zero at $v_0$: a point there with $f<0$ would allow an additional displacement increasing $\ell$ within $\Omega$. The real vertical Hessian of $f$ is thus nonnegative on $\ker\mathrm d\ell$. Consequently $\mathcal L_f$ is nonnegative on the complex space $V\cap T$, of dimension $r-1$. Hence $W\cap V=\{0\}$, and the projection $W\to T^{1,0}_{x_0}X$ is an isomorphism.

\emph{Step 3: plurisubharmonicity.} In local coordinates $z$ centered at $x_0$, choose a holomorphic section $\sigma_0$ through $v_0$ with tangent space $W$. The function $f\circ\sigma_0$ has zero first derivative and negative definite mixed Hessian. Its quadratic expansion is
\[
 f(\sigma_0(z))=\sum H_{\alpha\bar\beta}z_\alpha\bar z_\beta
                    +2\operatorname{Re}P(z)+O(|z|^3),
\]
where $P$ is a holomorphic quadratic polynomial. Choose a vertical vector $b$ with $\partial_vf(b)=1$, and replace $\sigma_0$ by $\sigma=\sigma_0-bP$. This preserves the first jet and cancels the pure quadratic part. Thus, after shrinking the chart,
\[
 f(\sigma(z))\leq-\epsilon|z|^2,
 \qquad \sigma(0)=v_0,
\]
for some $\epsilon>0$. In particular $\sigma(z)\in A_z$ throughout the chart.

Compactness of $\overline\Omega$ makes $p$ upper semicontinuous: for convergent covectors, the corresponding maximizing vectors have a subsequence whose limit lies in the limiting fiber of $\overline\Omega$. The supporting section gives the opposite inequality $p(z,\xi)\geq|\xi(\sigma(z))|$, with equality at $(x_0,\xi_0)$. Therefore $p$ is continuous off the zero section, and its upper bound gives continuity at zero. The holomorphic pairing $\xi(\sigma(z))$ is nonzero near the contact point, and its logarithmic absolute value is a pluriharmonic function lying below $\log p$ and equal to it at $(x_0,\xi_0)$. The mean-value property on small holomorphic disks therefore gives the submean inequality for $\log p$ at every point off the zero section. Thus $\log p$, and then $Q=\exp(2\log p)$ by Jensen's inequality, are plurisubharmonic there. At the zero section the continuous nonnegative function $Q$ vanishes, so its submean inequality on disks also holds there.

\emph{Step 4: uniform convexity in the fibers.} Consider all maximizing triples $(x,\xi,v)$ with $|\xi|_{k^*}=1$. They form a compact set, since $p$ is continuous, and by Step~2 the vertical differential $\mathrm d_vf$ does not vanish at these support points. Thus there are uniform constants $m>0$, $\rho>0$ and $C\geq1$ such that, for vertical vectors $w$,
\[
 |\mathrm d_vf|\geq m,\qquad
 f(v+w)\leq\mathrm d_vf(w)+C|w|_k^2\quad (|w|_k\leq\rho).
\]
Choose $0<\delta<\min\{a_0/2,\rho/2,m/(2C)\}$. At a support point, write $N$ for the outward real vertical unit normal and set $c=v-\delta N$. If $|w+\delta N|_k\leq\delta$, then
\[
 k_{\mathbb R}(w,N)\leq-\frac{|w|_k^2}{2\delta},\qquad
 f(v+w)\leq-\left(\frac{m}{2\delta}-C\right)|w|_k^2\leq0.
\]
All these points lie in the neighborhood where $f$ is defined, so the closed ball $\overline B_k(c,\delta)$ is contained in $A_x$.

The support function of this ball is $s_c(\xi)=|\xi(c)|+\delta|\xi|_{k^*}\leq p(x,\xi)$. Let $\xi$ be a covector of unit norm for which $v$ is a maximizing point, and let $\eta$ be $\xi$ with its phase normalized as in Step~2. Then $\operatorname{Re}\eta$ has real dual norm one, and its metric dual vector is $N$. Hence $\operatorname{Re}\eta(N)=1$, and since $|\eta(N)|\leq1$, also $\eta(N)=1$. We obtain $\eta(c)=a-\delta>0$, and therefore $s_c(\xi)=a$. Homogeneity now gives
\[
 Q(x,\xi)=\sup_c s_c(\xi)^2,
\]
where $c$ ranges over the centers of the supporting balls in the fiber. Every such center satisfies $|c|_k\leq b_0+\delta$, so by \cref{lem:ball-support-uniform-convexity} all the functions $s_c^2-\gamma|\xi|_{k^*}^2$ are convex for the same $\gamma>0$. Their supremum, which is $Q(x,\cdot)-\gamma|\cdot|_{k^*}^2$, is finite and therefore convex, as asserted.

\emph{Step 5: smoothing.} Finally, if $E$ is ample, \cref{prop:ample-dual-strict-homogeneous-function} gives the smooth strictly plurisubharmonic homogeneous function on $E^*$ required in \cref{prop:weak-homogeneous-convex-replacement}, and applying that proposition to $Q$ gives the smooth norm.
\end{proof}

For a complex manifold of dimension $N$ and an integer $1\leq q\leq N$, we say that the manifold is \emph{$q$-convex} if it has a smooth real exhaustion function with at least $N-q+1$ positive Levi eigenvalues outside a compact set; here an exhaustion function is one whose sublevel sets are compact. A $q$-convex manifold is also $q'$-convex for $q\leq q'\leq N$. We call such integers $q$ the \emph{smooth convexity indices} of the manifold, and the smallest of them, when one exists, its \emph{least smooth convexity index}.\index{$q$-convexity!least smooth convexity index}

\Needspace{9\baselineskip}
\begin{lemma}\label{lem:closed-submanifold-complement-index-bound}
\index{$q$-convexity!closed submanifold restriction}
Let $M$ be a complex manifold of dimension $N$, and let $S\subset M$ be a closed complex submanifold of dimension $D\geq1$. If $S$ has least smooth convexity index $\kappa$, then every smooth convexity index of $M$ is at least $\kappa$.
\end{lemma}
\begin{proof}
Suppose that $M$ is $p$-convex for an integer $1\leq p<\kappa$. Then a smooth exhaustion function $\psi$ has at least $N-p+1$ positive Levi eigenvalues outside a compact set $K\subset M$. Its restriction to the closed submanifold $S$ is an exhaustion function, since its sublevel sets are closed subsets of the compact sublevel sets of $\psi$, and $K\cap S$ is compact.

At a point of $S\setminus K$, choose a complex subspace $P\subset T^{1,0}M$ of dimension $N-p+1$ on which the Levi form of $\psi$ is positive definite. Then
\[
 \dim_{\mathbb C}(P\cap T^{1,0}S)\geq (N-p+1)+D-N=D-p+1.
\]
The Levi form of $\psi|_S$ is the restriction of the Levi form of $\psi$, since the inclusion is holomorphic, so it has at least $D-p+1$ positive eigenvalues. Since $p<\kappa\leq D$, this shows that $S$ is $p$-convex, which contradicts the definition of $\kappa$.
\end{proof}

In the following results, projectivizations parametrize complex lines.

\Needspace{9\baselineskip}
\begin{corollary}\label{cor:dual-norm-complement-convexity-obstruction}
\index{Projective bundle!complement convexity obstruction}
Suppose that $E$ is ample and that $E^*$ admits no smooth complex norm with plurisubharmonic square off the zero section. Then the complement
\[
 Z\setminus Y,\qquad Z=\mathbb P(\mathcal O_X\oplus E),
 \quad Y=\mathbb P(\mathcal O_X),
\]
is not $r$-convex.
\end{corollary}
\begin{proof}
The chart of $Z$ where the $\mathcal O_X$-coordinate is nonzero is identified with the total space of $E$, and $Y$ is its zero section. Suppose that $Z\setminus Y$ is $r$-convex. Since $\dim Z=n+r$, a smooth exhaustion function $\psi$ then has at least $n+1$ positive Levi eigenvalues outside a compact set $K$. Choose a relatively compact disk-bundle neighborhood $D\Subset E$ of $Y$. The set $Z\setminus D$ is compact and disjoint from $Y$. Choose a regular value $c$ of $\psi$ greater than its maximum on $Z\setminus D$ and on $K$.

Since $\{\psi\leq c\}$ is compact in $Z\setminus Y$, the set $\Omega=\{\psi>c\}\cup Y$ is an open neighborhood of $Y$. Its closure is a compact subset of $D$, and its boundary is the regular level set $\{\psi=c\}$. The defining function $f=c-\psi$ has at least $n+1$ negative Levi eigenvalues there. Restriction to the complex tangent hyperplane removes at most one negative direction, leaving at least $n$. Since $E$ is ample, \cref{thm:concave-neighbourhood-dual-norm} produces a smooth complex norm on $E^*$ with plurisubharmonic square, a contradiction.
\end{proof}

\Needspace{10\baselineskip}
\begin{corollary}\label{cor:hirzebruch-non-two-convex-complement}
\index{Hirzebruch surface!non-two-convex complements}
For every $e\geq0$, there is an ample rank-two bundle $E$ on the Hirzebruch surface $X=\mathbb F_e$ such that $\mathbb P(\mathcal O_X\oplus E)\setminus\mathbb P(\mathcal O_X)$ is not $2$-convex. The bundle may be chosen with the slope-stability and Chern classes of \cref{thm:hirzebruch-stable-ample-approximation}.
\end{corollary}
\begin{proof}
Use \cref{thm:hirzebruch-stable-ample-approximation} to choose $E$ whose dual Dolbeault operator lies in the $C^4$ neighborhood of \cref{thm:convex-norm-deformation-obstruction}; the split dual bundle satisfies the hypotheses of that theorem by \cref{prop:hirzebruch-positive-dual-tests}. Thus $E^*$ admits no smooth complex norm with plurisubharmonic square, and \cref{cor:dual-norm-complement-convexity-obstruction} with $r=2$ gives the assertion.
\end{proof}

\subsection{Curvature and projective complement exhaustions}
\label{subsec:projective-complement-curvature-index}

\Cref{cor:dual-norm-complement-convexity-obstruction} excludes $r$-convexity of the complement of the zero section in the projective completion, but it does not decide which larger indices occur. Upper bounds come from curvature. The reciprocal of the squared norm of a Hermitian metric is an exhaustion near the removed zero section, and every positive direction of the curvature, evaluated at a fiber vector, contributes a positive Levi eigenvalue. For a stable bundle of positive Gauduchon degree, the Hermitian--Einstein metric provides such a direction at every fiber vector.

In this subsection, $E\to X$ is a holomorphic vector bundle of rank $r\geq1$, $n\geq1$, and $h$ is a smooth Hermitian metric on $E$. In a holomorphic frame, write the matrix of $h$ as $M$, using column vectors and an inner product linear in its first argument, so that $h(v,v)=v^*Mv$. Put
\[
 F_h=\bar\partial(M^{-1}\partial M)
     =\sum_{\alpha,\beta}F_{\alpha\bar\beta}\,\mathrm dz_\alpha\wedge\mathrm d\bar z_\beta,
 \qquad R_h(u,\bar u)=\sum_{\alpha,\beta}F_{\alpha\bar\beta}u_\alpha\bar u_\beta.
\]
For $v\in E_x$, the real $(1,1)$-form $\alpha_{h,v}=\mathrm i\,h(F_hv,v)$ has the Hermitian coefficient form $u\mapsto h(R_h(u,\bar u)v,v)$. \index{Chern curvature!of a bundle metric}For a Hermitian form $\omega$ on $X$, the \emph{mean curvature} of $h$ is the $h$-Hermitian endomorphism $K_{\omega,h}=\Lambda_\omega(\mathrm iF_h)$, where $\Lambda_\omega$ denotes contraction with $\omega$, so that $\Lambda_\omega(\beta)\,\omega^n=n\beta\wedge\omega^{n-1}$ for $(1,1)$-forms $\beta$.\index{Mean curvature!of a bundle metric} The curvature $F_h$ is unscaled: for a line bundle, the curvature $-\ddc\log h$ of \cref{sec:convex-bundle-norm-obstruction} is $\mathrm iF_h/(2\pi)$.

Use the line convention for projectivization and put
\[
 U_E=\mathbb P(\mathcal O_X\oplus E)\setminus\mathbb P(\mathcal O_X).
\]
The chart $[1:v]$ identifies the total space of $E$ with an open subset of $\mathbb P(\mathcal O_X\oplus E)$, in which $\mathbb P(\mathcal O_X)$ is the zero section; thus $U_E$ contains $E^\times$. Its complex dimension is $n+r$. Convexity is understood in the sense of the smooth exhaustions of \cref{subsec:concave-neighbourhood-dual-norms}.

\Needspace{14\baselineskip}
\begin{lemma}\label{lem:reciprocal-bundle-norm-levi}
\index{Chern curvature!reciprocal norm Hessian}
Put $G(v)=h(v,v)$ and $\rho=G^{-1}$ on $E^\times$. In a holomorphic frame that is unitary and \emph{Chern-normal} at $x$, meaning that $M(x)$ is the identity and the first derivatives of $M$ vanish at $x$, the complex Hessian of $\rho$ at $(x,v)$ is
\[
 \mathcal L_\rho(u,w)
 =\frac{h(R_h(u,\bar u)v,v)}{G^2}
  -\frac{|w|_h^2}{G^2}
  +\frac{2|h(w,v)|^2}{G^3},
 \qquad u\in T_x^{1,0}X,\quad w\in E_x,
\]
where $(u,w)$ are the horizontal and vertical components of a tangent vector. There are no mixed horizontal--vertical terms, and the vertical block has exactly one positive and $r-1$ negative eigenvalues.
\end{lemma}
\begin{proof}
Such a frame is obtained by a constant change of frame making it unitary at $x$, followed by a holomorphic change of frame, linear in the coordinates, that cancels $\partial M$. Then $F_{\alpha\bar\beta}=-\partial_\alpha\partial_{\bar\beta}M$ at $x$. The horizontal first derivatives of $G$ and its mixed horizontal--vertical second derivatives vanish. The horizontal complex Hessian of $G$ is $-h(R_h(u,\bar u)v,v)$, its vertical complex Hessian is $|w|_h^2$, and its vertical first derivative in the direction $w$ is $h(w,v)$. Apply
\[
 \partial\bar\partial(G^{-1})
 =2G^{-3}\partial G\wedge\bar\partial G-G^{-2}\partial\bar\partial G.
\]
The vertical block equals $G^{-1}>0$ at $w=v$ and equals $-|w|_h^2/G^2$ when $w\perp_h v$. These two subspaces are orthogonal for the vertical block, which proves the statement about its eigenvalues.
\end{proof}

\Needspace{8\baselineskip}
\begin{theorem}\label{thm:bundle-curvature-complement-convexity}
\index{Projective bundle!curvature exhaustion}\index{$q$-convexity!bundle complements}
Let $s$ be an integer with $0\leq s\leq n$. Suppose that for every $x\in X$ and every $0\ne v\in E_x$, the coefficient form of $\alpha_{h,v}$ has at least $s$ positive eigenvalues. Then $U_E$ is $(n+r-s)$-convex.
\end{theorem}
\begin{proof}
By \cref{lem:reciprocal-bundle-norm-levi}, the horizontal block of $\mathcal L_\rho$ has at least $s$ positive eigenvalues and the vertical block has one. Their direct sum thus has at least $s+1$ positive eigenvalues at every point of $E^\times$.

Choose $\epsilon>0$ and a smooth function $\chi:[0,+\infty)\to[0,1]$ equal to one on $[0,\epsilon]$ and to zero on $[2\epsilon,+\infty)$. Define
\[
 \psi([1:v])=\frac{\chi(G(v))}{G(v)}\quad(v\ne0),
\]
and extend $\psi$ by zero on the divisor at infinity $\mathbb P(E)$. This extension is smooth. Indeed, near a point $[0:v]$ with $v\ne0$, use a local representative $[u:v]$ in a projective chart. On its intersection with the affine chart, $G(v/u)=h(v,v)/|u|^2$ tends to infinity as $u\to0$, so $\chi(G)$ vanishes in a neighborhood of the point.

The function $\psi$ tends to infinity near the removed zero section. Since the ambient projective bundle is compact, every sublevel set of $\psi$ is compact in $U_E$: it is closed there and cannot accumulate on the removed section. Put $D=\{G<\epsilon\}\subset E$, including the zero section, and $K=\mathbb P(\mathcal O_X\oplus E)\setminus D$. Then $K$ is a compact subset of $U_E$, and $\psi=\rho$ on $U_E\setminus K$. Thus the exhaustion $\psi$ has at least $s+1$ positive Levi eigenvalues outside a compact set. In dimension $n+r$, this is $(n+r-s)$-convexity.
\end{proof}

Taking $s=0$ gives $(n+r)$-convexity for every bundle. For $s=1$, the curvature hypothesis is precisely RC-positivity of the metric, as in \cref{def:directional-rc-positivity}. A positive definite mean curvature provides such a direction by taking a trace.

To apply this to stable bundles, we extend the degree of line bundles of \cref{sec:convex-bundle-norm-obstruction} to vector bundles through their determinants.

\Needspace{13\baselineskip}
\begin{definition}\label{def:gauduchon-bundle-slope}
\index{Gauduchon degree!vector bundles}\index{Slope stability!Gauduchon form}\index{Bott--Chern class!Gauduchon degree}
For a Gauduchon form $\omega$ on $X$, define
\[
 \deg_\omega E=\deg_{\omega^{n-1}}\det E
 =\int_X\frac{\mathrm i\operatorname{tr}F_h}{2\pi}\wedge\omega^{n-1},
 \qquad \mu_\omega(E)=\frac{\deg_\omega E}{r}.
\]
For a nonzero torsion-free coherent sheaf $\mathcal F$, use the determinant line bundle $\det\mathcal F=(\bigwedge^{\operatorname{rank}\mathcal F}\mathcal F)^{**}$ to define its degree and slope. The bundle $E$ is \emph{$\omega$-slope-stable} if $\mu_\omega(\mathcal F)<\mu_\omega(E)$ for every coherent subsheaf $\mathcal F\subset E$ with $0<\operatorname{rank}\mathcal F<r$ and torsion-free quotient.
\end{definition}

If $\det h'=\mathrm e^f\det h$, the two curvature representatives differ by $-\ddc f$, whose integral against $\omega^{n-1}$ vanishes by \cref{lem:ibp}. Thus the degree depends only on the Bott--Chern class $c_1^{\mathrm{BC}}(E)$ of the closed form $\mathrm i\operatorname{tr}F_h/(2\pi)$, and it is the pairing of this class with $\omega^{n-1}$. The Hopf example in \cref{ex:hopf-exact-complement-index} shows that the de Rham first Chern class of $E$ does not determine the degree.

\Needspace{10\baselineskip}
\begin{corollary}\label{cor:positive-mean-complement-convexity}
\index{Hermitian--Einstein metric!positive degree}\index{Mean curvature!complement convexity}
If $K_{\omega,h}$ is positive definite for a Hermitian form $\omega$ on $X$, then $U_E$ is $(n+r-1)$-convex. This holds, in particular, when $\omega$ is Gauduchon, $E$ is $\omega$-slope-stable, and $\deg_\omega E>0$.
\end{corollary}
\begin{proof}
For $v\ne0$, the trace of $\alpha_{h,v}$ with respect to $\omega$ is $h(K_{\omega,h}v,v)>0$, and a Hermitian form with positive trace has a positive eigenvalue. Apply \cref{thm:bundle-curvature-complement-convexity} with $s=1$.

Under the Gauduchon and stability hypotheses, the classical existence theorem for Hermitian--Einstein metrics on stable bundles gives a metric $h$ with $K_{\omega,h}=\lambda\operatorname{Id}_E$, where $\lambda$ is a constant. Taking the trace and integrating the contraction identity gives
\[
 r\lambda\int_X\omega^n
 =n\int_X\operatorname{tr}(\mathrm iF_h)\wedge\omega^{n-1}
 =2\pi n\deg_\omega E>0.
\]
Here $\operatorname{tr}F_h=\bar\partial\partial\log\det M$, so $\mathrm i\operatorname{tr}F_h/(2\pi)$ represents $c_1^{\mathrm{BC}}(E)$. Thus $\lambda>0$, and the first assertion applies.
\end{proof}

\Needspace{9\baselineskip}
\begin{corollary}\label{cor:hirzebruch-exact-complement-index}
\index{Hirzebruch surface!least complement convexity index}
For every $e\geq0$, there is a stable ample rank-two bundle $E$ on $X=\mathbb F_e$ as in \cref{cor:hirzebruch-non-two-convex-complement} for which $U_E$ is $3$-convex but not $2$-convex. Its least smooth convexity index is three.
\end{corollary}
\begin{proof}
Use the ruling section $S$ and fiber $f$ of \cref{lem:hirzebruch-ruling-sections}, and put $H=S+(3e+1)f$ and $P=4H-f$. Choose $E$ as in \cref{cor:hirzebruch-non-two-convex-complement}, $\mu_P$-stable and with $c_1(E)=4H$. The intersection numbers of \cref{lem:hirzebruch-ruling-sections} give
\[
 c_1(E)\cdot P=4H\cdot(4H-f)=80e+28>0.
\]
By the same lemma, $P=4S+(12e+3)f$ is very ample. Take for $\omega$ the pulled-back Fubini--Study K\"ahler form representing $c_1(\mathcal O_X(P))$. Degrees with respect to $\omega$ are intersection numbers with $P$, so $E$ is $\omega$-slope-stable with $\deg_\omega E>0$, and \cref{cor:positive-mean-complement-convexity} gives $3$-convexity, since $n=r=2$. By the choice of $E$, the complement is not $2$-convex. Finally, $1$-convexity would imply $2$-convexity, so the least index is three.
\end{proof}

\subsection{External tensor products and complement convexity}
\label{subsec:external-tensor-complement-curvature}

On a product of $k$ manifolds, the external tensor product of bundles with positive mean curvature has, at every nonzero fiber vector, a positive curvature direction tangent to each factor. Indeed, the curvature of the tensor product metric has no components mixing different factors, so this count applies to every fiber vector, including vectors that are not decomposable tensors. \Cref{thm:bundle-curvature-complement-convexity} then improves the convexity bound by $k$.

Let $k\geq2$. For $1\leq i\leq k$, let $X_i$ be a connected compact complex manifold of dimension $n_i\geq1$, let $E_i\to X_i$ be a holomorphic vector bundle of rank $r_i\geq1$, and choose a smooth Hermitian metric $h_i$ on $E_i$ and a Hermitian form $\omega_i$ on $X_i$. With $p_i:X=\prod_iX_i\to X_i$ the projections, put
\[
 E=\bigotimes_{i=1}^kp_i^*E_i,\qquad h=\bigotimes_{i=1}^kp_i^*h_i,
 \qquad n=\sum_i n_i,\quad r=\prod_i r_i.
\]
For an endomorphism $A_i$ of $E_i$, write $\widehat A_i$ for its action on the $i$th tensor factor, with the identity on every other factor. We use the Chern curvature, the forms $\alpha_{h,v}$ and the mean curvatures of \cref{subsec:projective-complement-curvature-index}.

\Needspace{14\baselineskip}
\begin{lemma}\label{lem:external-tensor-curvature-traces}
\index{External tensor product!curvature}
The curvature of $h$ is
\[
 F_h=\sum_i\widehat{p_i^*F_{h_i}}.
\]
For every $v\in E_x$, the coefficient form of $\alpha_{h,v}$ is block diagonal with respect to the decomposition $T_x^{1,0}X=\bigoplus_iT_{x_i}^{1,0}X_i$, and
\[
 \operatorname{tr}_{\omega_i}\bigl(\alpha_{h,v}|_{T_{x_i}X_i}\bigr)
 =h\bigl(\widehat K_i v,v\bigr),\qquad K_i=K_{\omega_i,h_i}.
\]
\end{lemma}
\begin{proof}
In product holomorphic frames, the matrix of the metric is $M=\bigotimes_iM_i$, and its connection matrix is
\[
 M^{-1}\partial M=\sum_i\widehat{M_i^{-1}\partial M_i}.
\]
Each summand depends only on the coordinates of $X_i$. Applying $\bar\partial$ gives the curvature formula and shows that all components mixing different factors vanish. Pairing with $v$ gives the block diagonal form. Contracting the $i$th block with $\omega_i$ commutes with the identity actions on the other tensor factors, which gives the trace identity.
\end{proof}

\Needspace{10\baselineskip}
\begin{theorem}\label{thm:external-tensor-complement-convexity}
\index{Projective bundle!external tensor complements}\index{$q$-convexity!external tensor products}
If every $K_i$ is positive definite, then
\[
 \mathbb P(\mathcal O_X\oplus E)\setminus\mathbb P(\mathcal O_X)
\]
is $(n+r-k)$-convex.
\end{theorem}
\begin{proof}
At each point, choose an $h_i$-orthonormal eigenbasis of $K_i$. The resulting product basis is $h$-orthonormal, and $\widehat K_i$ has the positive eigenvalues of $K_i$, repeated according to the dimensions of the other tensor factors. Hence $h(\widehat K_i v,v)>0$ for every nonzero $v$, whether or not $v$ is decomposable.

By \cref{lem:external-tensor-curvature-traces}, each diagonal block of $\alpha_{h,v}$ has positive trace and therefore a positive eigenvalue, so the block diagonal form has at least $k$ positive eigenvalues. Since $k\leq n$, \cref{thm:bundle-curvature-complement-convexity} with $s=k$ gives the assertion.
\end{proof}

\Needspace{12\baselineskip}
\begin{corollary}\label{cor:stable-factor-complement-convexity}
\index{Hermitian--Einstein metric!external tensor factors}
Suppose that each $\omega_i$ is Gauduchon, that $E_i$ is $\omega_i$-slope-stable as in \cref{def:gauduchon-bundle-slope}, and that $\deg_{\omega_i}E_i>0$. Then the complement $\mathbb P(\mathcal O_X\oplus E)\setminus\mathbb P(\mathcal O_X)$ is $(n+r-k)$-convex. For $k$ stable ample Hirzebruch examples as in \cref{cor:hirzebruch-exact-complement-index}, with their polarizations, this gives $(2^k+k)$-convexity.
\end{corollary}
\begin{proof}
As in the proof of \cref{cor:positive-mean-complement-convexity}, choose a Hermitian--Einstein metric on each $E_i$. Its mean curvature is $\lambda_i\operatorname{Id}_{E_i}$, with $\lambda_i>0$ by the positivity of the degree and the integrated trace identity. Apply \cref{thm:external-tensor-complement-convexity}. For the Hirzebruch factors, $n_i=r_i=2$, so $n+r-k=2k+2^k-k=2^k+k$.
\end{proof}

\subsection{Homogeneous plurisubharmonic functions and Hopf complements}
\label{subsec:flat-restriction-hopf-complements}

Let $F\to B$ be a holomorphic bundle of rank $r\geq1$ on a connected compact complex manifold of dimension $n\geq1$. Write $F^\times$ for its complement of the zero section, and put $N=n+r$ and
\[
 U_F=\mathbb P(\mathcal O_B\oplus F)\setminus\mathbb P(\mathcal O_B).
\]
The reciprocal-norm exhaustion in \cref{thm:bundle-curvature-complement-convexity} at $s=0$ makes $U_F$ $N$-convex. A homogeneous plurisubharmonic function forces this upper bound to be sharp. Flat restrictions then give a lower bound for the Hopf example.

\Needspace{10\baselineskip}
\begin{proposition}\label{prop:homogeneous-psh-sharp-complement-index}
\index{Plurisubharmonic function!homogeneous complement obstruction}\index{Projective bundle!least convexity index}
Suppose that $F^\times$ admits a smooth positive plurisubharmonic function $G$ satisfying $G(cv)=|c|^2G(v)$ for every $c\in\mathbb C^*$ and $v\in F^\times$. Then the least smooth convexity index of $U_F$ is $N$.
\end{proposition}
\begin{proof}
It remains to exclude $(N-1)$-convexity, since every smaller convexity index would imply it. Suppose that a smooth exhaustion $\psi$ has at least two positive Levi eigenvalues outside a compact set $K\subset U_F$.

Fix a reference Hermitian norm on $F$. On its compact unit sphere bundle, $G$ has a positive minimum and a finite maximum. Homogeneity consequently compares $G$ with that squared norm uniformly on $F^\times$. Thus
\[
 \Sigma_t=\{G=t^2\},\qquad t>0,
\]
are compact smooth hypersurfaces, and multiplication by $t$ identifies $\Sigma_1$ with $\Sigma_t$. Smoothness follows also from the Euler identity: if $R_v=v$ is the real radial vector at $v$, then $\mathrm dG(R_v)=2G(v)>0$. Since $K$ stays away from the removed zero section, there is $\epsilon>0$ such that every $\Sigma_t$, $0<t<\epsilon$, is disjoint from $K$.

Put
\[
 M(t)=\max_{v\in\Sigma_1}\psi([1:tv]),\qquad 0<t<\epsilon.
\]
On any closed interval inside $(0,\epsilon)$, the radial derivatives of the smooth function under the maximum are uniformly bounded over the compact sphere. Hence $M$ is locally Lipschitz. Properness of the exhaustion implies $M(t)\to+\infty$ as $t\to0$: otherwise maximizing points along a sequence tending to zero would lie in one compact sublevel while accumulating on the removed section. Choose $0<t_0<t_1<\epsilon$ with $M(t_0)>M(t_1)$. The absolute continuity of a Lipschitz function and its integral derivative formula give a differentiability point $t\in(t_0,t_1)$ with $M'(t)<0$.

Choose $v\in\Sigma_1$ attaining $M(t)$, and set $p=tv$. For both signs of $h$, the comparison
\[
 M(t+h)\geq\psi([1:(t+h)v])
\]
is an equality at $h=0$. Differentiability of $M$ therefore gives $\frac{\mathrm d}{\mathrm dt}\psi([1:tv])=M'(t)$. Since $\psi$ has a constrained maximum on $\Sigma_t$ at $p$ and $\mathrm dG\ne0$, there is a real $\lambda$ with $\mathrm d\psi=\lambda\,\mathrm dG$ at $p$. The radial derivative gives
\[
 \lambda=\frac{M'(t)}{2t}<0.
\]
The function $\psi-\lambda G$ has zero differential at $p$ and nonpositive real Hessian on $T_p\Sigma_t$. Indeed, its restriction to the sphere is $\psi$ minus a constant, and its zero differential removes the second derivative of a curve realizing a tangent vector. On the complex tangent hyperplane $H=T_p\Sigma_t\cap\mathrm iT_p\Sigma_t$, this implies
\[
 \mathcal L_\psi|_H\leq\lambda\mathcal L_G|_H\leq0.
\]
The first inequality follows by adding the real Hessian inequalities in a tangent vector and its complex rotation; the second uses plurisubharmonicity and $\lambda<0$. A two-dimensional complex subspace on which $\mathcal L_\psi$ is positive definite would meet the complex codimension-one space $H$ nontrivially, a contradiction. Thus $\psi$ has at most one positive Levi eigenvalue at $p\notin K$. Combine this contradiction with \cref{thm:bundle-curvature-complement-convexity} at $s=0$.
\end{proof}

\Needspace{10\baselineskip}
\begin{corollary}\label{cor:unitary-flat-exact-complement-index}\label{lem:unitary-flat-complement-obstruction}
\index{Unitary flat bundle!least complement convexity index}\index{Projective bundle!flat restriction}
If $F$ admits a Hermitian metric with flat Chern connection, then the least smooth convexity index of $U_F$ is $n+r$.
\end{corollary}
\begin{proof}
In parallel holomorphic unitary frames, the squared norm is $G(v)=\sum_j|v_j|^2$, which is plurisubharmonic on the total space. Constant unitary transitions make it global, smooth, positive on $F^\times$ and homogeneous of degree two under complex scalars. Apply \cref{prop:homogeneous-psh-sharp-complement-index}.
\end{proof}

\Needspace{8\baselineskip}
Fix $n\geq2$, a real $0<q<1$ and a complex $a\ne0$. On the scalar Hopf manifold
\[
 X_q=(\mathbb C^n\setminus\{0\})/\langle z\mapsto qz\rangle
\]
let $E_a$ be the holomorphic line bundle defined by $(z,v)\sim(qz,av)$. Put
\[
 b=-\frac{\log|a|}{\log q},
 \qquad G(z,v)=|z|^{2b}|v|^2.
\]
The identity $q^{2b}|a|^2=1$ makes $G$ the squared norm of a Hermitian metric on $E_a$.

\Needspace{10\baselineskip}
\begin{example}\label{ex:hopf-exact-complement-index}
\index{Hopf manifold!least complement convexity index}\index{Gauduchon degree!de Rham trivial line bundle}\index{$q$-convexity!Hopf complement}
The least smooth convexity index of $U_{E_a}$ is two when $0<|a|<1$, and $n+1$ when $|a|\geq1$. For every Gauduchon form $\gamma$ on $X_q$, the degree of $E_a$ is positive, zero or negative according as $|a|<1$, $|a|=1$ or $|a|>1$. Its de Rham first Chern class vanishes in all cases.
\end{example}
\begin{proof}
Write $\alpha=\ddc\log|z|^2$. This form descends because the contraction changes its potential by a constant. For $u\in T_z^{1,0}(\mathbb C^n\setminus\{0\})$, its coefficient form, apart from the positive factor $1/(2\pi)$, is
\[
 \frac{|u|^2}{|z|^2}-\frac{|\langle u,z\rangle|^2}{|z|^4},
\]
where $\langle\cdot,\cdot\rangle$ is the Euclidean Hermitian inner product. Cauchy--Schwarz shows that it is positive semidefinite with kernel $\mathbb Cz$, hence of rank $n-1$. The normalized line curvature of $G$ is $-b\alpha$. The integral of $\alpha\wedge\gamma^{n-1}$ is strictly positive, so the degree has the sign of $-b$, giving the three stated cases. The map $z\mapsto(z/|z|,\log|z|)$ identifies the covering space with $S^{2n-1}\times\mathbb R$, and the contraction translates the second coordinate by $\log q$. Hence $X_q$ is diffeomorphic to $S^{2n-1}\times S^1$ and $H^2_{\mathrm{dR}}(X_q,\mathbb R)=0$.

If $|a|<1$, then $b<0$, so every nonzero fibre vector supplies $n-1$ positive curvature directions. Apply \cref{thm:bundle-curvature-complement-convexity} with $r=1$ and $s=n-1$ to obtain $2$-convexity. The map $\pi:X_q\to\mathbb P^{n-1}$, $[z]\mapsto[z]$, has compact embedded elliptic fibres $B\simeq\mathbb C^*/q^{\mathbb Z}$. On a lifted fibre $z=\zeta z_0$, the logarithm of the metric coefficient is $b\log|\zeta|^2+b\log|z_0|^2$, which is pluriharmonic. Thus $E_a|_B$ has flat Chern connection. The complement $U_{E_a|_B}$ is a closed complex submanifold of $U_{E_a}$. A $1$-convex exhaustion of the latter would restrict to a proper exhaustion of the former, strictly plurisubharmonic outside the compact intersection of its exceptional set with the submanifold. This contradicts \cref{cor:unitary-flat-exact-complement-index}, since the restricted bundle has base dimension one and rank one. The phase of $a$ affects the gluing but not these metric and curvature calculations.

If $|a|\geq1$, then $b\geq0$. On the covering of $E_a^\times$, the function
\[
 \log G=b\log|z|^2+\log|v|^2
\]
is plurisubharmonic: the first summand has semipositive Levi form and the second is pluriharmonic for $v\ne0$. Consequently
\[
 \mathcal L_G=G\bigl(\mathcal L_{\log G}+\partial\log G\otimes\overline{\partial\log G}\bigr)\geq0.
\]
The descended $G$ is smooth, positive and homogeneous of degree two under complex scalars. Apply \cref{prop:homogeneous-psh-sharp-complement-index} with rank one and total dimension $n+1$.
\end{proof}

